%%%%%%%%%%%%%%%%%%%%%%%%%%%%%%%%%%%%%%%%%%%%%%%%%%%%%%%%%%%%%%%%%%%%%%
% Lớp 10 | Xác suất | Xác suất của biến cố
% Dạng 0D0?2-5: Liên quan việc chọn đối tượng khác
% Số câu: 77  (sinh tự động bởi tach_cau.py; từ nay thầy sửa trực tiếp file này)
%%%%%%%%%%%%%%%%%%%%%%%%%%%%%%%%%%%%%%%%%%%%%%%%%%%%%%%%%%%%%%%%%%%%%%

%%%%%%%%%%%%%%%% TRẮC NGHIỆM 4 PHƯƠNG ÁN %%%%%%%%%%%%%%%%

% nguồn: MaHoaNH_TracNghiem_Lop10
\begin{ex}%[0D0N2-5]%[ma:c7a95642]
	Rút ra một lá bài từ bộ bài $52$ lá. Xác suất để  được lá bích là
	\choice 
		{$\dfrac{1}{13}$} 
		{\True $\dfrac{1}{4}$} 
		{$\dfrac{12}{13}$} 
		 {$\dfrac{3}{4}$}
	\loigiai{Số phần tử của không gian  mẫu là $n\left(\Omega\right)=52$.\\
		Gọi $A$ là biến cố: \lq  \lq Rút được lá bích \rq \rq. Khi đó $n(A)=13$.\\
	Vậy  $\mathrm{P} \left(A\right)=\dfrac{n(A)}{n(\Omega)}=\dfrac{13}{52}=\dfrac{1}{4}$.}
\end{ex}

% nguồn: MaHoaNH_TracNghiem_Lop10
\begin{ex}%[0D0N2-5]%[ma:0a17c11e]
	Gieo một con súc sắc cân đối và đồng  chất. Xác suất xuất hiện mặt có số chấm chia hết cho 3 bằng
	\choice 
		{$\dfrac{1}{3}$} 
		{\True  $\dfrac{1}{3}$}
		 {$\dfrac{2}{3}$}
		 {$\dfrac{2}{3}$}
	 \loigiai{
		Số phần tử của không gian  mẫu là $n(\Omega)=6$. \\Gọi $A$ là biến cố \lq\lq xuất hiện mặt có số chấm chia hết cho 3\rq\rq. \\
	Số phần  tử của biến cố $A$ là $n(A)=2$. \\Vậy xác suất của biến cố $A$ là $\mathrm{P}(A)=\dfrac{1}{3}$.}
\end{ex}

% nguồn: MaHoaNH_TracNghiem_Lop10
\begin{ex}%[0D0N2-5]%[ma:28f83419]
	Một tổ có $6$ nam và $4$ nữ. Chọn  ngẫu nhiên hai người. Tính xác suất sao cho trong hai người được chọn có ít nhất một người là nữ.
	 \choice
		{$\dfrac{4}{5}$} 
		 {\True $\dfrac{2}{3}$}
		{$\dfrac{2}{15}$} 
		{$\dfrac{1}{3}$} 
	 \loigiai{
		Số phần tử không gian mẫu là số cách  chọn 2 người trong 10 người $n(\Omega)= \mathrm{C}_{10}^2$.\\
		Gọi $A$ là biến cố hai người có ít nhất một người  là nữ.\\
	Suy ra $\overline{A}$ là biến cố \lq\lq cả hai người là nam\rq\rq. Ta có $n\left(\overline{A}\right)=\mathrm{C}_6^2$. Do  đó xác suất $\mathrm{P}(A)=1-\dfrac{\mathrm{C}_6^2}{\mathrm{C}_{10}^2}=\dfrac{2}{3}$.}
\end{ex}

% nguồn: MaHoaNH_TracNghiem_Lop10
\begin{ex}%[0D0N2-5]%[ma:c07c8ea7]
	Lấy ngẫu nhiên $3$ bi từ một hộp chứa $4$ bi đỏ và $5$ bi vàng. Tính xác suất để $3$ bi lấy ra có ít nhất một bi đỏ. 
	\choice 
		{$\dfrac{3}{4}$} 
		 {$\dfrac{10}{21}$}
		{$\dfrac{2}{7}$} 
		{\True $\dfrac{37}{42}$} 
	\loigiai{ 
		Số phần tử của không gian mẫu là  $n(\Omega)=\mathrm{C}_9^3=84$.\\
		Số phần tử của biến cố $A$ là $n(A)=\mathrm{C}_4^3+\mathrm{C}_4^2  \mathrm{C}_5^1+\mathrm{C}_4^1 \mathrm{C}_5^2=74$.
	\\Xác suất  để 3 bi được chọn có ít nhất 1 bi đỏ là $\mathrm{P}(A)=\dfrac{n(A)}{n(\Omega)}=\dfrac{74}{84}=\dfrac{37}{42}$.}
\end{ex}

% nguồn: MaHoaNH_TracNghiem_Lop10
\begin{ex}%[0D0N2-5]%[ma:99026fb3]
	Từ các  chữ số $1; 2; 4; 6; 8; 9$ lấy ngẫu nhiên mộtt số. Xác suất để lấy được số lẻ bằng
	\choice 
		{\True $\dfrac{1}{3}$} 
		{$\dfrac{1}{2}$} 
		 {$\dfrac{1}{4}$}
		 {$\dfrac{1}{6}$}
	 \loigiai{
		Số phần tử của không  gian mẫu là $n(\Omega)=\mathrm{C}_6^1=6$.\\
		Gọi $A$ là biến  cố \lq\lq lấy được số lẻ\rq\rq. Số phần tử của biến cố là $n(A)=\mathrm{C}_2^1=2$.\\
	Xác suất của biến cố $A$ là $\mathrm{P}(A)=\dfrac{n(A)}{n(\Omega)}=\dfrac{1}{3}$.} 
\end{ex}

% nguồn: MaHoaNH_TracNghiem_Lop10
\begin{ex}%[0D0N2-5]%[ma:5da090cd]
	Trong trò chơi \lq\lq Hãy chọn giá đúng\rq\rq\,chiếc kim của bánh xe có thể dừng lại ở 1 trong 20 nấc điểm với khả năng như nhau. Tính xác xuất để trong hai lần quay, chiếc kim của  bánh xe đó dừng lại ở hai nấc điểm khác nhau.
	 \choice
		 {$\dfrac{1}{20}$}
		 {\True $\dfrac{19}{20}$}
		 {$\dfrac{1}{10}$}
		{$\dfrac{9}{10}$} 
	\loigiai{ 
		 Số phần tử của không gian mẫu là $n(\Omega)=20^2=400$.\\
		Gọi $A$ là biến cố \lq\lq trong hai lần quay, chiếc kim của bánh xe đó dừng lại ở  hai nấc điểm khác nhau\rq\rq. \\
	Số phần tử của biến cố $A$ là $n(A)=\mathrm{A}_{20}^2=380$.\\Vậy xác suất của biến cố $A$ là  $\mathrm{P}(A)=\dfrac{n(A)}{n(\Omega)}=\dfrac{380}{400}=\dfrac{19}{20}$.}
\end{ex}

% nguồn: MaHoaNH_TracNghiem_Lop10
\begin{ex}%[0D0N2-5]%[ma:b99739f8]
	Một chiếc hộp chứa $20$ viên bi, trong đó có $8$ viên bi màu đỏ, $7$  viên bi màu xanh và $5$ viên bi màu vàng. Lấy ngẫu nhiên ra $3$ viên bi. Xác suất để $3$ viên bi lấy ra đều màu đỏ là
	\choice 
		{\True  $\dfrac{14}{285}$}
		{$\dfrac{45}{128}$} 
		 {$\dfrac{28}{145}$}
		 {$\dfrac{15}{248}$}
	\loigiai{ Số cách lấy ngẫu nhiên ra $3$ viên bi từ $20$ viên bi trong hộp là $n \left(\Omega\right)= \mathrm{C}_{20}^3=1140$.\\ 
		Số cách  chọn ra $3$ viên bi màu đỏ là $n (A)=\mathrm{C}_8^3=56$.\\
		Xác  suất để $3$ viên bi lấy ra đều màu đỏ là $P(A)=\dfrac{56}{1140}=\dfrac{14}{285}$.
	 }
\end{ex}

% nguồn: MaHoaNH_TracNghiem_Lop10
\begin{ex}%[0D0N2-5]%[ma:99c5b286]
	Một hộp chứa  $3$ quả cầu màu trắng và $2$ quả cầu màu đen. Chọn ngẫu nhiên đồng thời từ trong hộp ra $2$ quả. Xác suất chọn được hai quả màu trắng bằng
	\choice 
		 {$\dfrac{1}{3}$}
		 {$\dfrac{1}{5}$}
		 {\True $\dfrac{3}{10}$}
		{$\dfrac{2}{5}$} 
	\loigiai{ Chọn  ngẫu nhiên đồng thời từ trong hộp ra $2$ quả. Số phần tử của không gian mẫu là $$n(\Omega)=\mathrm{C}_{5}^2=10.$$
		Số cách chọn từ trong hộp ra $2$ quả màu trắng là $n(A)=\mathrm{C}_3^2=3$.\\ 
		Xác suất để chọn được  hai quả màu trắng là $P(A)=\dfrac{3}{10}$.
	} 
\end{ex}

% nguồn: MaHoaNH_TracNghiem_Lop10
\begin{ex}%[0D0N2-5]%[ma:e499f4c4]
	Từ một hộp đựng $4$ cái bút bi và $5$ cái bút chì, lấy ngẫu nhiên hai cái bút. Xác suất để lấy được cả hai cái bút bi  là
	\choice 
		{$\dfrac{1}{3}$} 
		 {$\dfrac{4}{9}$}
		 {$\dfrac{5}{6}$}
		 {\True $\dfrac{1}{6}$}
	 \loigiai{
		Lấy ngẫu nhiên $2$ cái bút trong hộ  gồm $9$ cây bút có $\mathrm{C}_{9}^{2}$ cách. Do đó $n(\Omega)=\mathrm{C}_{9}^{2}$.\\
		 Gọi $A$ là biến cố: \lq \lq Lấy được cả hai cái bút bi\rq \rq.\\
		Lấy $2$ trong $4$ bút bi có $\mathrm{C}_4^2$ cách nên $n(A)=\mathrm{C}_4^2$.\\ 
	 Do đó $\mathrm{P}(A)=\dfrac{n(A)}{n(\Omega)}=\dfrac{\mathrm{C}_4^2}{\mathrm{C}_{9}^{2}}=\dfrac{1}{6}$.}
\end{ex}

% nguồn: MaHoaNH_TracNghiem_Lop10
\begin{ex}%[0D0N2-5]%[ma:242217e1]
	Gieo một đồng xu cân đối đồng chất liên tiếp hai  lần. Tính xác suất để cả hai lần gieo đều được mặt sấp.
	\choice 
		{\True $\dfrac{1}{4}$} 
		 {$\dfrac{1}{6}$}
		 {$\dfrac{1}{8}$}
		 {$\dfrac{1}{2}$}
	 \loigiai{
		Số phần tử của  không gian mẫu là $n(\Omega)=2\cdot2=4$.\\
		Gọi $A$ là  biến cố \lq\lq cả hai lần gieo đều mặt sấp\rq\rq.
		\\Số phần tử của biến cố $A$ là  $n(A)=1$.
	Vậy xác suất  của biến cố $A$ là $\mathrm{P}(A)=\dfrac{n(A)}{n(\Omega)}=\frac{1}{4}$.}
\end{ex}

% nguồn: MaHoaNH_TracNghiem_Lop10
\begin{ex}%[0D0N2-5]%[ma:fef8b3ae]
	Một hộp  đựng $5$ viên bi đỏ, $4$ viên bi xanh. Lấy ngẫu nhiên $3$ viên bi từ hộp đó. Tính xác suất lấy được ít nhất một viên đỏ.
	\choice 
		 {$\dfrac{1}{21}$}
		{$\dfrac{37}{42}$} 
		 {$\dfrac{5}{42}$}
		 {\True $\dfrac{20}{21}$}
	 \loigiai{
		Số phần tử của không gian  mẫu là $n(\Omega)=\mathrm{C}_9^3=84$.\\
		Gọi  $A$ là biến cố cần tìm. \\
		Số phần tử của biến cố  $A$ là $n(A)= \mathrm{C}_5^3+\mathrm{C}_5^2 \cdot \mathrm{C}_4^1+\mathrm{C}_5^1 \cdot \mathrm{C}_4^2=80$.\\
	Vậy xác suất của biến  cố $A$ là $\mathrm{P}(A)=\dfrac{80}{84}=\dfrac{20}{21}$.}
\end{ex}

% nguồn: MaHoaNH_TracNghiem_Lop10
\begin{ex}%[0D0N2-5]%[ma:008d056f]
	Gieo lần lượt 2 con xúc xắc 6 mặt. Xác  suất để tổng số chấm xuất hiện trên hai mặt của hai con súc sắc bằng $7$ là
	 \choice
		 {$\dfrac{6}{7}$}
		{$\dfrac{1}{7}$} 
		{\True$\dfrac{1}{6}$} 
		 {$\dfrac{5}{6}$}
	\loigiai{ 
		Số phần tử của không  gian mẫu là $n(\Omega)=6 \cdot 6=36$.\\
		Gọi $A$ là biến cố \lq\lq Tổng số chấm xuất hiện trên hai mặt của hai con súc  sắc bằng 7\rq\rq.\\
		Ta có $A=\lbrace(1; 6),(2;  5),(3; 4),(4; 3),(5; 2),(6; 1)\rbrace \Rightarrow n(A)=6$.\\
	Vậy xác suất của biến cố $A$  là $\mathrm{P}(A)=\dfrac{6}{36}=\dfrac{1}{6}$.}
\end{ex}

% nguồn: MaHoaNH_TracNghiem_Lop10
\begin{ex}%[0D0N2-5]%[ma:fba2ff19]
	Một đoàn đại biểu gồm $5$ người được chọn ra  từ một tổ gồm $8$ nam và $7$ nữ để tham dự hội nghị. Xác suất để chọn được đoàn đại biểu có đúng $2$ người nữ là
	\choice 
		{\True $\dfrac{56}{143}$} 
		 {$\dfrac{140}{429}$}
		{$\dfrac{1}{143}$} 
		 {$\dfrac{28}{715}$}
	 \loigiai{
		Số phần  tử của không gian mẫu là $n(\Omega)=\mathrm{C}_{15}^5$.\\
		Gọi $A$  là biến cố \lq\lq Chọn được đoàn đại biểu có đúng 2 người nữ\rq\rq. \\
		Số  phần tử của biến cố $A$ là $n(A)=\mathrm{C}_7^2 \cdot \mathrm{C}_8^3 $.\\
	 Vậy xác suất cần tìm là $\mathrm{P}(A)=\dfrac{n(A)}{n(\Omega)}=\dfrac{56}{143}$.}
\end{ex}

% nguồn: MaHoaNH_TracNghiem_Lop10
\begin{ex}%[0D0N2-5]%[ma:5738e501]
	Từ một hộp chứa $11$ quả cầu màu đỏ và $4$ quả cầu màu xanh, lấy ngẫu nhiên đồng thời $3$ quả cầu. Xác suất để lấy được $3$ quả cầu màu xanh bằng 
	\choice 
		{$\dfrac{4}{165}$} 
		{\True $\dfrac{4}{455}$} 
		 {$\dfrac{33}{91}$}
		{$\dfrac{24}{455}$} 
	\loigiai{ 
		Số  phần tử của không gian mẫu là $n(\Omega)=\mathrm{C}_{15}^3=455$.\\
		Gọi $A$ là biến cố \lq\lq lấy được 3 quả cầu  màu xanh\rq\rq.\\ Số phần tử của biến cố $A$ là $n(A)=\mathrm{C}_4^3=4$.\\
		Vây xác suất của biến cố $A$ là $\mathrm{P}(A)=\dfrac{4}{455}$. 
	} 
\end{ex}

% nguồn: MaHoaNH_TracNghiem_Lop10
\begin{ex}%[0D0N2-5]%[ma:24d113d2]
	Một hội nghị có $15$ nam và $6$ nữ. Chọn ngẫu nhiên $3$ người vào ban  tổ chức. Xác suất để $3$ người lấy ra là nam
	\choice 
		 {$\dfrac{1}{2}$}
		{\True  $\dfrac{91}{266}$}
		{$\dfrac{4}{33}$} 
		 {$\dfrac{1}{11}$}
	 \loigiai{
		 Gọi ${A}$ là biến cố: \lq\lq 3 người lấy ra là nam\rq\rq.\\
		Số phần tử của không  gian mẫu là $n(\Omega)=\mathrm{C}_{21}^3=1\,330$.
		\\Số phần tử của biến cố $A$ là  $n(A)=\mathrm{C}_{15}^3=455$.\\
		Vậy xác suất  để 3 người lấy ra là nam là $\mathrm{P}(A)=\dfrac{n(A)}{n(\Omega)}=\dfrac{13}{38}=\dfrac{91}{266}$.
	} 
\end{ex}

% nguồn: MaHoaNH_TracNghiem_Lop10
\begin{ex}%[0D0N2-5]%[ma:2deaebcb]
	Tìm trên giá sách có $ 4$ quyển sách Toán, $3$ quyển sách Vật lý, $2$ quyển sách Hoá học. Lấy ngẫu nhiên $3$ quyển sách trên kệ sách đó. Tính xác suất để $3$ quyển được lấy ra  đều là sách Toán.
	\choice 
		 {$\dfrac{2}{7}$}
		 {\True $\dfrac{1}{21}$}
		 {$\dfrac{37}{42}$}
		 {$\dfrac{5}{42}$}
	\loigiai{ 
		Số phần tử của không  gian mẫu là $n\left(\Omega\right)=\mathrm{C}_9^3=84$.\\
		Gọi $A$ là biến cố \lq\lq $3$ quyển được lấy ra đều  là sách Toán\rq\rq.\\
		Số kết quả thuận lợi của biến cố $A$ là  $n\left(A\right)=\mathrm{C}_4^3=4$.\\
		Xác suất biến  cố $A$ là $ \mathrm{P}(A)=\dfrac{1}{21}$.
	 }
\end{ex}

% nguồn: MaHoaNH_TracNghiem_Lop10
\begin{ex}%[0D0H2-5]%[ma:cc3de764]
	Quỳnh có một hộp hình lập phương mà sáu mặt được tô bởi ba màu đỏ, vàng, xanh sao cho chỉ có các mặt đối nhau thì tô cùng màu. Hỏi tung ngẫu nhiên hộp đó thì xác suất để được mặt xanh ngửa là bao  nhiêu?
	\choice 
		{\True  $\dfrac{1}{3}$}
		 {$\dfrac{2}{3}$}
		 {$\dfrac{1}{2}$}
		 {$\dfrac{1}{6}$}
	\loigiai{Do không gian mẫu của phép thử có $6$ kết quả, trong đó có $2$ kết quả thuận lợi cho biến cố đang xét (do có hai mặt màu xanh) nên xác suất của biến cố: \lq \lq  Thu được mặt xanh ngửa\rq \rq là $\dfrac{2}{6}=\dfrac{1}{3}$.}
\end{ex}

% nguồn: MaHoaNH_TracNghiem_Lop10
\begin{ex}%[0D0H2-5]%[ma:1c80569f]
	Một hộp chứa $5$ bi xanh, $4$ bi đỏ. Chọn  ngẫu nhiên $2$ bi từ hộp này. Xác suất để chọn được $2$ bi cùng màu là
	 \choice
		{$\dfrac{2}{9}$} 
		{$\dfrac{5}{9}$} 
		 {$\dfrac{1}{9}$}
		{\True $\dfrac {4}{9}$} 
	\loigiai{Số phần tử của không  gian mẫu $n(\Omega)=\mathrm{C}_{9}^2$.\\
		Gọi $A$ là biến cố \lq\lq chọn được $2$  bi cùng màu\rq\rq.\\
		Ta có  $n(A)=\mathrm{C}_{5}^2+ \mathrm{C}_{4}^2$.\\
	Vậy xác suất cần tìm là  $$P(A)=\dfrac{n(A)}{n(\Omega)}=\dfrac{\mathrm{C}_{5}^2+ \mathrm{C}_{4}^2}{\mathrm{C}_{9}^2}=\dfrac {4}{9}.$$}
\end{ex}

% nguồn: MaHoaNH_TracNghiem_Lop10
\begin{ex}%[0D0H2-5]%[ma:5b96f422]
	Trên giá sách có $5$ quyển sách Toán, $4$ quyển sách Lý và $3$ quyển sách Hóa. Lấy ngẫu nhiên $3$ quyển sách. Xác suất để $3$ quyển sách được lấy  ra thuộc $3$ môn khác nhau là
	\choice 
		 {$\dfrac{8}{11}$}
		{\True $\dfrac{3}{11}$} 
		 {$\dfrac{1}{110}$}
		 {$\dfrac{109}{110}$}
	\loigiai{ 
		Số cách lấy ngẫu nhiên ba quyển sách  trên giá là $n(\Omega)=\mathrm{C}_{12}^{3}$.\\
		Gọi $A$ là biến  cố: \lq \lq Lấy được $3$ quyển sách thuộc $3$ môn khác nhau\rq \rq.\\
		Khi đó ta cần chọn lần lượt $1$ quyển sách Toán, $1$ quyển sách Lý, $1$ quyển sách Hóa nên ${n(A)}=\mathrm{C}_{5}^{1} \cdot  \mathrm{C}_{4}^{1} \cdot \mathrm{C}_{3}^{1}$.\\
	Xác suất để lấy được $3$ quyển sách thuộc $3$ môn khác nhau  là $\mathrm{P}(A)=\dfrac{\mathrm{C}_{5}^{1} \cdot \mathrm{C}_{4}^{1} \cdot \mathrm{C}_{3}^{1}}{\mathrm{C}_{12}^{3}}=\dfrac{3}{11}$.}
\end{ex}

% nguồn: MaHoaNH_TracNghiem_Lop10
\begin{ex}%[0D0H2-5]%[ma:88220bec]
	Từ một hộp chứa $10$ quả bóng gồm $4$ quả màu đỏ và $6$ quả màu  xanh, lấy ngẫu nhiên đồng thời $3$ quả. Xác suất của biến cố \lq\lq lấy được $3$ quả màu xanh\rq\rq\ là
	 \choice
		{\True  $\dfrac{\mathrm{C}_6^3}{\mathrm{C}_{10}^3}$}
		 {$\dfrac{\mathrm{C}_4^3}{\mathrm{C}_{10}^3}$}
		 {$\dfrac{\mathrm{C}_{10}^3}{\mathrm{C}_{10}^4}$}
		{$\dfrac{\mathrm{C}_4^3}{\mathrm{C}_6^3}$} 
	 \loigiai{
		Số phần tử của không gian mẫu là  số cách lấy $3$ quả bóng trong $10$ quả bóng, do đó $n(\Omega)=\mathrm{C}_{10}^3$.\\
		Số phần tử của biến cố \lq\lq lấy được $3$ quả  màu xanh\rq\rq là $\mathrm{C}_6^3$.\\
		Suy  ra xác suất là $\dfrac{\mathrm{C}_6^3}{\mathrm{C}_{10}^3}$.
	} 
\end{ex}

% nguồn: MaHoaNH_TracNghiem_Lop10
\begin{ex}%[0D0H2-5]%[ma:dc23cf35]
	Một túi chứa $6$ bi  xanh, $4$ bi đỏ. Lấy ngẫu nhiên $2$ bi. Tính xác suất để
	được cả hai bi  đều màu đỏ.
	 \choice
		 {$\dfrac{5}{12}$}
		{\True  $\dfrac{2}{15}$}
		{$\dfrac{7}{45}$} 
		 {$\dfrac{8}{15}$}
	\loigiai{ 
		Gọi $A$ là  biến cố lấy được cả hai viên bi màu đỏ ta có $n(A)=\mathrm{C}^1_4\cdot \mathrm{C}^1_3=12$
		và không  gian mẫu $n(\Omega)=\mathrm{C}^2_{10} = 90.$ \\
		Vậy xác suất cần tìm là $\mathrm{P}(A)=\dfrac{12}{90}=\dfrac{2}{15}.$ 
	} 
\end{ex}

% nguồn: MaHoaNH_TracNghiem_Lop10
\begin{ex}%[0D0H2-5]%[ma:669face3]
	Từ một hộp chứa $12$ quả bóng gồm $5$ quả màu đỏ và $7$ quả màu  xanh, lấy ngẫu nhiên đồng thời $3$ quả. Xác suất để lấy được $3$ quả màu xanh bằng
	 \choice
		 {$\dfrac{5}{12}$}
		 {$\dfrac{2}{7}$}
		 {$\dfrac{1}{22}$}
		{\True  $\dfrac{7}{44}$}
	 \loigiai{
		Số phần tử của không gian  mẫu là $n(\Omega)=\mathrm{C}_{12}^{3}$.\\
		Gọi $A$ là biến cố: \lq\lq $3$ quả lấy  được đều có màu xanh\rq\rq.\\
		Ta có  $n(A)=\mathrm{C}_{7}^{3}$.\\
		Xác suất của  biến cố $A$ là $\mathrm{P}(A)=\dfrac{n(A)}{n(\Omega)}=\dfrac{\mathrm{C}_{7}^{3}}{\mathrm{C}_{12}^{3}}=\dfrac{7}{44}$.
	 }
\end{ex}

% nguồn: MaHoaNH_TracNghiem_Lop10
\begin{ex}%[0D0H2-5]%[ma:52f644c5]
	Có hai hộp bút chì màu. Hộp thứ nhất có $5$ bút chì màu đỏ khác nhau và $7$ bút chì màu xanh khác nhau. Hộp thứ hai có $8$ bút chì màu đỏ khác nhau và $4$ bút chì màu xanh khác nhau. Chọn ngẫu nhiên hai cây bút  chì. Số cách để có $1$ cây bút chì màu đỏ và $1$ cây bút chì màu xanh là
	 \choice
		{\True $143$} 
		{$35$} 
		{$32$} 
		{$24$} 
	 \loigiai{
		\begin{itemize}
			\item Số cách để chọn được $1$ cây bút chì màu đỏ  là $5+8=13$ cách.
			\item Số cách để  chọn được $1$ cây bút chì màu xanh là $7+4=11$ cách.
		\end{itemize}
		Vậy  có $13\cdot11=143$ cách.
	 }
\end{ex}

% nguồn: MaHoaNH_TracNghiem_Lop10
\begin{ex}%[0D0H2-5]%[ma:2c7c1487]
	Một lô hàng gồm $1000$ sản phẩm,  trong đó có $50$ phế phẩm. Lấy ngẫu nhiên từ lô hàng đó $1$ sản phẩm. Xác suất để lấy được sản phẩm tốt là
	 \choice
		{$0{,}94$} 
		{$0{,}96$} 
		{\True  $0{,}95$}
		 {$0{,}97$}
	\loigiai{ 
		Số phần  tử của không gian mẫu là $n(\Omega) = \mathrm{C}_{1000}^{1}$.\\
		 Gọi $A\colon$ \lq\lq Lấy được sản phẩm tốt \rq\rq.\\
		 Khi đó $n(A) = \mathrm{C}_{950}^{1}$.\\
		Xác suất để lấy được sản phẩm tốt là 
		$$ \mathrm{P}(A) = \dfrac{n(A)}{n(\Omega)} =  \dfrac{\mathrm{C}_{950}^{1}}{\mathrm{C}_{1000}^{1}} = \dfrac{19}{20} = 0{,}95. $$
	} 
\end{ex}

% nguồn: MaHoaNH_TracNghiem_Lop10
\begin{ex}%[0D0H2-5]%[ma:6e871352]
	Gieo một con súc sắc cân đối đồng chất ba lần.  Xác suất tích số chấm trong ba lần gieo bằng $6$ là
	 \choice
		{$\dfrac{1}{2}$} 
		{$\dfrac{5}{108}$} 
		 {$\dfrac{5}{9}$}
		{\True $\dfrac{1}{24}$} 
	\loigiai{ 
		 Ta có $n(\Omega)=6\times 6\times 6=216$.\\
		Gọi $A$ là biến cố \lq\lq tích  số chấm trong ba lần gieo bằng $6$\rq\rq.\\
		$A=\left\{(1,2,3);(1,3,2);(2,1,3);(2,3,1);(3,1,2);(3,2,1);(1,1,6);(1,6,1);(6,1,1)\right\}$.\\ 
		Suy ra,  $n(A)=9$.\\
		Vậy  xác suất của biến cố $A$ là $\mathrm{P}(A)=\dfrac{n(A)}{n(\Omega)}=\dfrac{9}{216}=\dfrac{1}{24}$.
	} 
\end{ex}

% nguồn: MaHoaNH_TracNghiem_Lop10
\begin{ex}%[0D0H2-5]%[ma:9dbde3b4]
	Trong một chiếc hộp có $3$  quả cầu xanh và $2$ quả cầu đỏ (các quả cầu này đều khác nhau), rút ngẫu nhiên hai quả cầu từ hộp này, xác suất để hai quả cầu lấy được luôn có quả cầu đỏ là
	 \choice
		 {\True $\dfrac{7}{10}$}
		 {$\dfrac{7}{20}$}
		 {$\dfrac{6}{10}$}
		 {$\dfrac{8}{10}$}
	\loigiai{ 
		\begin{itemize}
			\item Số  phần tử của không gian mẫu là $n(\Omega)=\mathrm{C}_5^2=10$.
			\item Gọi $A$  là biến cố:\,\lq\lq Hai quả cầu lấy được luôn có quả cầu đỏ\rq\rq.\\
			Số phần  tử của biến cố $A$ là $n(A)=\mathrm{n}_3^1\cdot\mathrm{C}_2^1+\mathrm{C}_2^2=7$.\\
			Xác  suất của biến cố $A$ là $\mathrm{P}(A)=\dfrac{\mathrm{n}(A)}{\mathrm{n}(\Omega)}=\dfrac{7}{10}$.
		\end{itemize}
	 }
\end{ex}

% nguồn: MaHoaNH_TracNghiem_Lop10
\begin{ex}%[0D0H2-5]%[ma:132eef0f]
	Từ  một hộp chứa $5$ quả cầu trắng và $3$ quả cầu đen, lấy ngẫu nhiên đồng thời $2$ quả. Tính xác suất để hai quả đó cùng màu.
	\choice 
		{$ \dfrac{5}{14} $} 
		{\True $ \dfrac{13}{28} $} 
		{$  \dfrac{15}{28} $}
		{$ \dfrac{3}{28}  $}
	\loigiai{ 
		Ta có $  n(\Omega)=\mathrm{C}_8^2=28 $.\\
		Gọi $ A $ là biến cố \lq\lq Hai quả lấy ra cùng  màu\rq\rq.
		\begin{itemize}
			\item Trường hợp 1: Hai  quả cầu cùng màu trắng có $ \mathrm{C}_5^2=10 $.
			\item Trường hợp 2: Hai quả cầu cùng màu đen có $ \mathrm{C}_3^2=3  $.
		\end{itemize}
		Do đó $ n(A)=10+3=13  $.\\
		Vậy xác suất của biến cố  $ A $ là $ \mathrm{P}(A)=\dfrac{n(A)}{n(\Omega)}=\dfrac{13}{28} $.
	} 
\end{ex}

% nguồn: MaHoaNH_TracNghiem_Lop10
\begin{ex}%[0D0H2-5]%[ma:b8d420d2]
	Một  hộp chứa $15$ viên bi kích thước như nhau, trong đó có $6$ viên bi màu trắng được đánh số từ $1$ đến $6$, có $5$ viên bi màu đỏ được đánh số từ $1$ đến $5$ và $4$ viên bi màu vàng được đánh số từ $1$ đến $4$. Lấy ngẫu nhiên $2$ viên bi từ hộp. Tính xác suất để $2$ viên bi được lấy vừa khác màu vừa khác số.
	\choice 
		 {$\dfrac{57}{105}$}
		{$\dfrac{61}{107}$} 
		{$\dfrac{15}{26}$} 
		{\True $\dfrac{61}{105}$} 
	\loigiai{ 
		Số phần tử của không gian mẫu  là $n\left(\Omega\right)=\mathrm{C}_{15}^2=105$.\\
		Gọi  $A$ là biến cố \lq\lq $2$ bi lấy ra vừa khác màu vừa khác số \rq\rq.
		\begin{itemize}
			\item {\bf Trường hợp 1:} $2$ bi lấy ra  có màu trắng và đỏ.\\
			Lấy $1$ bi đỏ có: $\mathrm{C}_5^1$  cách.\\
			 Lấy $1$ bi trắng khác số với bi đỏ có $\mathrm{C}_5^1$ cách.\\
			Suy ra trường hợp  này có $\mathrm{C}_5^1\cdot\mathrm{C}_5^1=25$ cách.
			\item {\bf Trường hợp 2:} $2$  bi lấy ra có màu vàng và đỏ.\\
			Lấy $1$ bi vàng có: $C_4^1$  cách.\\
			 Lấy $1$ bi đỏ khác số với bi vàng có $\mathrm{C}_4^1$ cách.\\
			Suy ra trường hợp này có $\mathrm{C}_4^1\cdot\mathrm{C}_4^1=16$  cách.
			\item {\bf Trường  hợp 3:} $2$ bi lấy ra có mày trắng và vàng.\\
			Lấy $1$ bi  vàng có: $\mathrm{C}_4^1$ cách.\\
			Lấy  $1$ bi trắng khác số với bi vàng có $\mathrm{C}_5^1$ cách.\\
			Suy ra  trường hợp này có $\mathrm{C}_4^1\cdot \mathrm{C}_5^1=20$ cách.
		\end{itemize}
		Do đó $n(A)=25+16+20=61$ cách chọn $2$ viên bi được lấy vừa khác màu vừa khác số.\\ 
		Suy  ra $\mathrm{P}(A)=\dfrac{n(A)}{n\left(\Omega\right)}=\dfrac{61}{105}$.
	 }
\end{ex}

% nguồn: Ok_TracNghiemLop11_Lan1
\begin{ex}%[0D0H2-5]%[ma:25a3cef8]
	Tại một hội thảo quốc tế có $50$ nhà khoa học, trong đó có $31$ người thành thạo tiếng Anh, $21$ người thành thạo tiếng Pháp và $5$ người thành thạo cả tiếng Anh và tiếng Pháp. Chọn ngẫu nhiên một người trong hội thảo. Xác suất để người được chọn không thông thạo cả hai thứ tiếng Anh hoặc Pháp  là
	\choice 
		 {$\dfrac{7}{50}$}
		 {\True$\dfrac{3}{50}$}
		 {$\dfrac{9}{50}$}
		{$\dfrac{11}{50}$} 
	\loigiai{ 
		Gọi biến cố $C$ \lq\lq Người được chọn không thông thạo cả hai thứ tiếng Anh hoặc Pháp\rq\rq. Ta thấy, $C$ và $A\cup B$  là hai biến cố đối nhau. Do đó,
	$$P(C)=1-P(A\cup  B)=1-\frac{47}{50}=\frac{3}{50}.$$}
\end{ex}

% nguồn: Ok_TracNghiemLop11_Lan1
\begin{ex}%[0D0H2-5]%[ma:52327f02]
	Tại một hội thảo quốc tế có $50$ nhà khoa học, trong đó có $31$ người thành thạo tiếng Anh, $21$ người thành thạo tiếng Pháp và $5$ người thành thạo cả tiếng Anh và tiếng Pháp. Chọn ngẫu nhiên một người trong hội thảo. Xác suất để người được chọn thành  thạo ít nhất một trong hai thứ tiếng Anh hoặc Pháp là
	 \choice
		 {\True$\dfrac{47}{50}$}
		 {$\dfrac{37}{50}$}
		 {$\dfrac{39}{50}$}
		 {$\dfrac{41}{50}$}
	 \loigiai{
		Gọi biến cố $A$ \lq\lq Người được chọn thành  thạo tiếng An\rq\rq.\\
		Biến cố $B$ \lq\lq Người được chọn thành  thạo tiếng Pháp\rq\rq.\\
		Biến cố $A\cup B$ \lq\lq Người  được chọn thông thạo tiếng Anh hoặc tiếng Pháp\rq\rq\\
		Biến cố $AB$ \lq\lq Người được chọn thông thạo cả tiếng Anh và tiếng  Pháp\rq\rq.
		Khi đó 
		$$P(A)=\frac{31}{50}; \  P(B)=\frac{21}{50};\ P(AB)=\frac{5}{50}.$$
	Do đó, $P(A\cup  B)=P(A)+P(B)-P(AB)=\dfrac{31}{50}+\dfrac{21}{50}-\dfrac{5}{50}=\dfrac{47}{50}$.}
\end{ex}

% nguồn: Ok_TracNghiemLop11_Lan1
\begin{ex}%[0D0H2-5]%[ma:b041540b]
	Một hộp chứa 15 quả cầu gồm 4 quả cầu màu xanh, 5 quả cầu màu đỏ  và 6 quả cầu vàng. Chọn ngẫu nhiên 4 quả cầu từ hộp đó. Xác suất để 4 quả cầu được chọn ra cùng màu bằng
	 \choice
		{$\dfrac{5}{91}$} 
		 {$\dfrac{7}{15}$}
		 {$\dfrac{5}{11}$}
		{\True $\dfrac{1}{65}$} 
	 \loigiai{
		Số phần tử không gian mẫu $n\left(  \Omega \right)=\mathrm{C}_{15}^4=1365$. \\
		Gọi  $A$ là biến cố: “lấy được $4$ quả cầu cùng màu”. \\
		 Khi đó, $n(A)=\mathrm{C}_4^4+\mathrm{C}_5^4+\mathrm{C}_6^4=21$. \\
		Xác suất $\mathrm{P}(A)=\dfrac{n(A)}{n\left( \Omega \right)}=\dfrac{1}{65}$. 
	 }
\end{ex}

% nguồn: MaHoaNH_TracNghiem_Lop10
\begin{ex}%[0D0V2-5]%[ma:46981f40]
	Để kiểm tra chất lượng sản phẩm từ công ty sữa, người ta gửi đến bộ phận kiểm nghiệm $5$ hộp sữa cam,  $4$ hộp sữa dâu và $3$ hộp sữa nho. Bộ phận kiểm nghiệm chọn ngẫu nhiên $3$ hộp để phân tích mẫu. Xác suất để $3$ hộp sữa được chọn có cả $3$ loại là
	 \choice
		{\True  $\dfrac{3}{11}$}
		{$\dfrac{1}{110}$} 
		{$\dfrac{3}{55}$} 
		 {$\dfrac{1}{22}$}
	\loigiai{ 
		$n(\Omega)=\mathrm{C}_{12}^3=220$.\\ 
		Gọi $A$ là biến cố \lq\lq $3$ hộp  sữa được chọn có cả $3$ loại\rq\rq.
		\begin{itemize}
			\item Chọn $1$ hộp sữa  cam có $5$ cách.
			\item Chọn $1$ hộp sữa dâu có $4$ cách. 
			\item  Chọn $1$ hộp sữa nho có $3$ cách.
		\end{itemize}
		Suy ra $n(A)=5\cdot  4\cdot 3=60$.\\
		Vậy xác suất  $\mathrm{P}(A)=\dfrac{n(A)}{n(\Omega)}=\dfrac{60}{220}=\dfrac{3}{11}$.
	} 
\end{ex}

% nguồn: MaHoaNH_TracNghiem_Lop10
\begin{ex}%[0D0V2-5]%[ma:70bb69aa]
	Chọn ngẫu nhiên $13$ lá bài từ bộ bài tú lơ khơ gồm $52$ lá. Tính xác suất để $13$ lá bài được chọn có  $6$ lá ghép được thành \lq\lq ba đôi thông\rq\rq, nghĩa là ba đôi có thứ tự liên tiếp nhau (không tính đôi $2$).
	 \choice
		{\True  $\dfrac{7128}{39151}$}
		 {$\dfrac{7129}{39151}$}
		 {$\dfrac{7130}{39151}$}
		{$\dfrac{7131}{39151}$} 
	\loigiai{ 
		Số phần tử của không  gian mẫu $n(\Omega) = \mathrm{C}_{52}^{13}$.\\
		Gọi $A$ là biến cố \lq\lq  6 lá ghép được thành ba đôi thông\rq\rq.\\
		Ta  có
		\begin{itemize}
			\item Có $10$ cách chọn ba giá trị  liên tiếp từ $3$ đến $A$ (bỏ đôi $2$).
			\item Với mỗi giá trị, có $\mathrm{C}_4^2$ cách chọn đôi (trong $4$ lá  mỗi giá trị).
			\item Còn lại chọn $7$ lá từ  $46$ lá còn lại $\mathrm{C}_{46}^7$.
		\end{itemize}
		$\Rightarrow n(A) = 10 \cdot \mathrm{C}_4^2  \cdot \mathrm{C}_4^2 \cdot \mathrm{C}_4^2 \cdot \mathrm{C}_{46}^7 = 7128$.\\
		Vậy xác suất là  $P(A) = \dfrac{7128}{39151}$.
	 }
\end{ex}


%%%%%%%%%%%%%%%% ĐÚNG/SAI %%%%%%%%%%%%%%%%

% nguồn: MaHoaNH_DungSai_Lop10
\begin{ex}%[0D0H2-5]%[ma:fe2d8556]
	Gieo một con súc sắc cân đối và đồng chất. Giả  sử súc sắc xuất hiện mặt $b$ chấm. Tính xác suất để phương trình $x^2+bx+2=0$ có hai nghiệm phân biệt.
	\choiceTF 
		{$\dfrac{1}{2}$}{$\dfrac{5}{6}$}{$\dfrac{1}{3}$}{\True $\dfrac{2}{3}$} 
	 \loigiai{
		Để phương trình có hai nghiệm phân biệt thì $b^2-8>0\Leftrightarrow \hoac{& b>2\sqrt{2} \\  & b<-2\sqrt{2}.}$\\
		 Vì $b$ nguyên dương nên $b\in\{3;4;5;6\}$.\\
		Vậy xác suất cần tìm là $\mathrm{P}(A)=\dfrac{4}{6}=\dfrac{2}{3}$. 
	 }
\end{ex}

% nguồn: MaHoaNH_DungSai_Lop10
\begin{ex}%[0D0H2-5]%[ma:a99b737f]
	Một hộp chứa $5$ quả cầu xanh, $4$ quả cầu đỏ và $4$ quả cầu  vàng, lấy ngẫu nhiên ra $4$ quả, tính xác suất để lấy được $4$ quả có đủ $3$ màu.
	\choiceTF 
		{$\dfrac{2}{61347}$} 
		{$\dfrac{64}{143}$} 
		 {\True $\dfrac{80}{143}$}
		 {$\dfrac{76}{143}$}
	 \loigiai{
		Số phần tử không  gian mẫu $n(\Omega) = \mathrm{C}^4_{13} = 715$.\\
		Gọi biến cố $A$ là \rq\rq Lấy được $4$ quả có  đủ $3$ màu\lq\lq\ .
		Khi đó $n(A) = \mathrm{C}^2_{5}\mathrm{C}^1_{4}\mathrm{C}^1_{4} + \mathrm{C}^1_{5}\mathrm{C}^2_{4}\mathrm{C}^1_{4} + \mathrm{C}^1_{5}\mathrm{C}^1_{4}\mathrm{C}^2_{4}  = 400$.\\
		Xác suất của biến cố $A$ là $\mathrm{P}(A) = \dfrac{n(A)}{n(\Omega)}  = \dfrac{400}{715} = \dfrac{80}{143}$.
	} 
\end{ex}

% nguồn: MaHoaNH_DungSai_Lop10
\begin{ex}%[0D0H2-5]%[ma:22e6c427]
	Bộ bài tú lơ khơ có $52$ quân bài. Rút ngẫu nhiên  ra $4$ quân bài. Khi đó
	\choiceTF 
		{\True Số phần tử của không gian  mẫu là $270\,725$}
		{\True Xác suất của biến cố $B$: \lq\lq  Rút ra được tứ quý K\rq\rq\ là $\dfrac{1}{270\,725}$}
		{Xác suất  của biến cố $C$: \lq\lq $4$ quân bài rút ra có ít nhất một con Át\rq\rq\ là $\dfrac{1\,229}{54\,145}$}
		{\True Xác suất của biến cố  $D$: \lq\lq $4$ quân bài lấy ra có ít nhất hai quân bích\rq\rq\ là $\dfrac{5\,359}{20\,825}$}
	\loigiai{ 
		\begin{itemchoice}
			\itemch Ta có số cách chọn ngẫu nhiên  $4$ quân bài là: $\mathrm{C}_{52}^4=270\,725$ suy ra $n(\Omega)=270\,725$.
			\itemch Trong bộ bài chỉ có $1$ tứ quý $K$ nên ta  có $n(B)=1$\\
			 Vậy $\mathrm{P}(B)=\dfrac{1}{270\,725}$.
			 \itemch Ta có $\mathrm{C}_{48}^4$ cách rút $4$ quân bài mà không có con Át nào, suy ra $4$ quân bài rút ra có ít nhất một con Át là $n(C)=\mathrm{C}_{52}^4-\mathrm{C}_{48}^4\Rightarrow \mathrm{P}(C)=\dfrac{15\,229}{54\,145}$.
			\itemch Trong bộ bài có $13$ quân bích, số cách rút ra bốn quân bài mà trong đó số quân  bích không ít hơn hai là \[\mathrm{C}_{13}^2\cdot \mathrm{C}_{39}^2+\mathrm{C}_{13}^3\mathrm{C}_{39}^1+\mathrm{C}_{13}^4\cdot \mathrm{C}_{39}^0=69\,667.\]
			Suy ra $n(D)=69\,667\Rightarrow  \mathrm{P}(D)=\dfrac{5\,359}{20\,825}$.
		\end{itemchoice}
	} 
\end{ex}

% nguồn: MaHoaNH_DungSai_Lop10
\begin{ex}%[0D0H2-5]%[ma:1345dc93]
	Cho một thùng gồm $4$ bi xanh, $5$ bi đỏ và  $6$ bi vàng. Khi đó
	 \choiceTF
		{\True Tổng số bi trong thùng là $15$} 
		{\True Số  kết quả có thể khi lấy ra $2$ viên bi từ số bi trong thùng là $\mathrm{C}_{15}^2=105$}
		{Số kết quả lấy ra hai bi khác màu  từ số bi trong thùng là $76$}
		{\True Gọi $A$ là biến cố \lq\lq Lấy ra hai viên bi khác màu\rq\rq\  Xác suất của biến cố $A$ là $\mathrm{P}(A)=\dfrac{74}{105}$}
	\loigiai{\begin{itemchoice}
		\itemch Tổng số bi  trong thùng là $4+5+6=15$ (bi).\\
		Vậy tổng  số bi trong thùng là $15$.
		\itemch Số kết quả  có thể khi lấy ra $2$ viên bi bất kì từ $15$ viên bi là $\mathrm{C}_{15}^2=105$.\\
		 Vậy số kết quả có thể khi lấy ra $2$ viên bi từ số bi trong thùng là $\mathrm{C}_{15}^2=105$.
		\itemch Số kết  quả thuận lợi khi lấy ra hai bi khác màu là $\mathrm{C}_4^1\cdot \mathrm{C}_5^1+\mathrm{C}_5^1\cdot \mathrm{C}_6^1+\mathrm{C}_4^1\cdot \mathrm{C}_6^1=74$.\\
		Vậy số kết quả lấy ra hai bi khác màu từ  số bi trong thùng là $74$.
		\itemch Gọi $A$ là biến cố lấy ra hai viên bi khác  màu.\\
		Xác suất xảy ra $A$ là  $\mathrm{P}(A)=\dfrac{74}{105} \approx 70{,}5\%$.
	\end{itemchoice}
} 
\end{ex}

% nguồn: MaHoaNH_DungSai_Lop10
\begin{ex}%[0D0H2-5]%[ma:e28a41d6]
	Lấy ngẫu nhiên hai thẻ từ một chiếc hộp chứa  $20$ thẻ được đánh số từ $1$ đến $20$. Khi đó
	 \choiceTF
		{\True Số phần tử của không gian mẫu  là $190$}
		{\True Số phần tử  của biến cố lấy được hai thẻ mang số lẻ là $45$}
		 {Xác suất để hai thẻ lấy ra có tổng chia hết cho $2$ là $\dfrac{9}{38}$}
		{\True Xác suất để hai thẻ lấy ra có tích chia hết  cho $2$ là $\dfrac{29}{38}$}
	\loigiai{\begin{itemchoice}
		\itemch  Số phần tử của không gian mẫu là $n(\Omega)=\mathrm{C}_{20}^2=190$.
		\itemch Số phần tử của biến cố \lq\lq Lấy được hai thẻ mang  số lẻ\rq\rq\ là $\mathrm{C}_{10}^2=45$.
		\itemch Chọn hai thẻ mang số chẵn $\mathrm{C}_{10}^2$.\\ 
		 Chọn hai thẻ mang số lẻ $\mathrm{C}_{10}^2$.\\
		Suy ra số phần tử của biến cố \lq\lq Hai thẻ lấy ra có  tổng chia hết cho $2$\rq\rq\ là $\mathrm{C}_{10}^2+\mathrm{C}_{10}^2=90$.\\
		Xác suất của biến cố \lq\lq Hai thẻ lấy ra  có tổng chia hết cho $2$\rq\rq\ là $\dfrac{90}{190}=\dfrac{9}{19}$.
		\itemch  Chọn hai thẻ mang số chẵn $\mathrm{C}_{10}^2$.\\
		Chọn một thẻ  mang số chẵn và một thẻ mang số lẻ $\mathrm{C}_{10}^1\cdot \mathrm{C}_{10}^1$.\\
		Suy ra số phần tử của biến cố \lq\lq Hai thẻ lấy ra có tích chia hết cho $2$\rq\rq\  là $\mathrm{C}_{10}^2+\mathrm{C}_{10}^1\cdot \mathrm{C}_{10}^1=145$.\\
		Xác suất của biến cố  \lq\lq Hai thẻ lấy ra có tích chia hết cho $2$\rq\rq\ là $\dfrac{145}{190}=\dfrac{29}{38}$.
	\end{itemchoice}
} 
\end{ex}

% nguồn: MaHoaNH_DungSai_Lop10
\begin{ex}%[0D0H2-5]%[ma:d78ba9e8]
	Có $10$ thẻ được đánh số $1, 2,\ldots, 10$. Bốc ngẫu nhiên $2$ thẻ. Các khẳng  định sau đúng hay sai.
	\choiceTF 
		{\True Xác suất của biến cố tích $2$ số ghi trên $2$ thẻ bốc được là số  lẻ là $\dfrac{2}{9}$}
		{Xác suất của  biến cố chọn được $2$ thẻ mang số chẵn là $\dfrac{7}{9}$}
		{\True Xác suất của biến cố tổng $2$ số ghi trên $2$ thẻ bốc được là số  chẵn là $\dfrac{4}{9}$}
		{Xác suất của biến cố chọn được  $2$ thẻ mang số nguyên tố là $\dfrac{5}{9}$}
	\loigiai{ 
		\begin{itemchoice}
			\itemch Số  phần tử của không gian mẫu là $n(\Omega)=\mathrm{C}_{10}^2$.\\
			Gọi $ A $ là biến cố \lq\lq tích $2$ số ghi trên $2$ thẻ bốc được là  một số lẻ\rq\rq $\Rightarrow n(A)=\mathrm{C}_5^2$.\\
			Vậy  $\mathrm{P}(A)=\dfrac{n(A)}{n(\Omega)}=\dfrac{2}{9}$.
			\itemch Số phần tử của không gian mẫu  là $n(\Omega)=\mathrm{C}_{10}^2$.\\
			Gọi $ A $ là biến cố \lq\lq tổng $2$ số ghi trên $2$ thẻ bốc  được là một số chẵn\rq\rq $\Rightarrow n(A)=\mathrm{C}_5^2+\mathrm{C}_5^2$.\\
			Vậy  $\mathrm{P}(A)=\dfrac{n(A)}{n(\Omega)}=\dfrac{2}{9}$.
			\itemch Số phần tử của không gian mẫu  là $n(\Omega)=\mathrm{C}_{10}^2$.\\
			Gọi $ A $ là biến cố \lq\lq  chọn được $2$ thẻ mang số chẵn\rq\rq $\Rightarrow n(A)=\mathrm{C}_5^2$.\\
			Vậy $\mathrm{P}(A)=\dfrac{n(A)}{n(\Omega)}=\dfrac{4}{9}$. 
			\itemch Số phần  tử của không gian mẫu là $n(\Omega)=\mathrm{C}_{10}^2$.\\
			Gọi $ A  $ là biến cố \lq\lq chọn được $2$ thẻ mang số nguyên tố\rq\rq $\Rightarrow n(A)=\mathrm{C}_5^2$.\\
			 Vậy $\mathrm{P}(A)=\dfrac{n(A)}{n(\Omega)}=\dfrac{2}{15}$.
		\end{itemchoice}
	} 
\end{ex}

% nguồn: MaHoaNH_DungSai_Lop10
\begin{ex}%[0D0H2-5]%[ma:133230e8]
	Một cừa hàng kem có $8$ vị: trà xanh, sô cô la, cà phê, dâu, chanh, va ni, đào, chuối.  An muốn mua một cốc kem. Các khẳng định sau đúng hay sai.
	 \choiceTF
		{ Xác suất của biến cố cốc kem có đúng 2  vị là $\dfrac{5}{51}$}
		{\True Xác suất của  biến cố cốc kem có ít nhất 2 vị là $\dfrac{247}{255}$}
		{\True Xác suất của  biến cố cốc kem có nhiều nhất 4 vị là $\dfrac{54}{85}$}
		{Xác suất của biến cố cốc kem có 3 vị  trong đó có vị sô cô la là $\dfrac{8}{85}$}
	\loigiai{ 
		\begin{itemchoice}
			\itemch Số phần tử của không gian  mẫu là $n(\Omega)=\displaystyle\sum\limits_{i=1}^8 \mathrm{C}_8^i=255$.\\
			Gọi $ A $ là biến cố \lq\lq cốc kem  có đúng 2 vị\rq\rq $\Rightarrow n(A)=\mathrm{C}_8^2=28$.\\
			Vậy $\mathrm{P}(A)=\dfrac{n(A)}{n(\Omega)}=\dfrac{28}{255}$. 
			\itemch Số phần tử của không gian mẫu là $n(\Omega)=\displaystyle\sum\limits_{i=1}^8  \mathrm{C}_8^i=255$.\\
			Gọi $ A $ là biến cố \lq\lq cốc kem có đúng 2  vị\rq\rq $\Rightarrow n(A)=\mathrm{C}_8^2+\mathrm{C}_8^3+\mathrm{C}_8^4+\mathrm{C}_8^5+\mathrm{C}_8^6+\mathrm{C}_8^7+\mathrm{C}_8^8=247$.\\
			Vậy $\mathrm{P}(A)=\dfrac{n(A)}{n(\Omega)}=\dfrac{247}{255}$. 
			\itemch Số  phần tử của không gian mẫu là $n(\Omega)=\displaystyle\sum\limits_{i=1}^8 \mathrm{C}_8^i=255$.\\
			Gọi $ A $ là biến cố \lq\lq cốc kem có  nhiều nhất 4 vị\rq\rq $\Rightarrow n(A)=\mathrm{C}_8^1+\mathrm{C}_8^2+\mathrm{C}_8^3+\mathrm{C}_8^4=162$.\\
			 Vậy $\mathrm{P}(A)=\dfrac{n(A)}{n(\Omega)}=\dfrac{54}{85}$.
			\itemch Số phần  tử của không gian mẫu là $n(\Omega)=\displaystyle\sum\limits_{i=1}^8 \mathrm{C}_8^i=255$.\\
			Gọi $ A $ là biến cố \lq\lq cốc kem có 3 vị trong đó có vị  sô cô la\rq\rq $\Rightarrow n(A)=\mathrm{C}_7^2=21=21$.\\
			Vậy  $\mathrm{P}(A)=\dfrac{n(A)}{n(\Omega)}=\dfrac{7}{85}$.
		\end{itemchoice}
		d) Sai. An mua một cốc kem có 3 vị trong đó có vị sô cô la có  số cách chọn là $C_7^2=21$
	} 
\end{ex}

% nguồn: MaHoaNH_DungSai_Lop10
\begin{ex}%[0D0H2-5]%[ma:6ae69ae6]
	Một hộp chứa $10$ tấm thẻ được đánh số từ $1$ đến  $10$. Rút ngẫu nhiên hai tấm thể thẻ từ hộp đó. Khi đó
	 \choiceTF
		{Số phần tử của  không gian mẫu là $90$}
		{\True Xác suất để rút được hai tấm thẻ được  đánh số cùng chia hết cho $2$ là $\dfrac{2}{9}$}
		{Xác suất để  rút được hai tấm thẻ được đánh số đều là số nguyên tố là $\dfrac{1}{15}$}
		{\True Xác suất để rút được hai tấm thẻ có tổng là một lẻ là $\dfrac{5}{9}$} 
	\loigiai{ 
		\begin{itemchoice}
			 \itemch Chọn $2$ số ngẫu nhiên từ $1$ đến $10$, suy ra $n(\Omega)=\mathrm{C}_{10}^2=45$.
			\itemch Xét biến cố $A$: \lq\lq Rút được cả hai tấm thẻ được đánh số cùng chia hết  cho $2$\rq\rq.\\
			Ta có các số chia hết cho $2$ và không lớn hơn $10$ là $2$, $4$, $6$, $8$, $10$.\\ 
			Chọn ngẫu nhiên $2$ số trong $5$ số trên. Suy ra  $n(A)=\mathrm{C}_5^2$.\\
			Do vậy $\mathrm{P}(A)=\dfrac{\mathrm{C}_5^2}{\mathrm{C}_{10}^2}=\dfrac{2}{9}$. 
			\itemch Xét biến cố $B$: \lq\lq Rút được  hai tấm thẻ được đánh số đều là số nguyên tố\rq\rq.\\
			Ta có các số nguyên tố và không lớn  hơn $10$ là $2$, $3$, $5$, $7$.\\
			Chọn ngẫu nhiên $2$ số trong $4$  số trên. Suy ra $n(B)=\mathrm{C}_4^2$.\\
			Do vậy $\mathrm{P}(B)=\dfrac{\mathrm{C}_4^2}{\mathrm{C}_{10}^2}=\dfrac{2}{15}$. 
			\itemch Xét biến cố $D$: \lq\lq Rút được hai tấm thẻ có tổng là một  lẻ\rq\rq.\\
			Suy ra trong hai tấm thẻ phải có $1$  thẻ mang số lẻ và $1$ thẻ mang số chẵn.\\
			Chọn  $1$ thẻ mang số lẻ, có $5$ cách chọn.\\
			Chọn  $1$ thẻ mang số chẵn, có $5$ cách chọn.\\
			 Suy ra $n(D)=5\cdot 5=25$. \\
			Do vậy $\mathrm{P}(B)=\dfrac{25}{\mathrm{C}_{10}^2}=\dfrac{5}{9}$. 
		\end{itemchoice}
	 }
\end{ex}

% nguồn: MaHoaNH_DungSai_Lop10
\begin{ex}%[0D0H2-5]%[ma:f149f9d4]
	Có $100$ tấm thẻ được đánh số từ $1$ đến $100$. Lấy ngẫu  nhiên $5$ thẻ. Khi đó
	\choiceTF 
		{\True Số phần tử của không gian mẫu  là $\mathrm{C}_{100}^5$}
		{Xác suất để $5$ thẻ lấy ra đều  mang số chẵn là $0{,}5$}
		{\True Xác suất để $5$ thẻ lấy ra có $2$ thẻ mang số chẵn và $3$ thẻ mang số lẻ xấp  xỉ bằng $ 0{,}32$}
		{Xác suất để có ít nhất một số ghi  trên thẻ được chọn chia hết cho $3$ xấp xỉ bằng $0{,}78$}
	\loigiai{ 
		\begin{itemchoice}
			\itemch 
			Vì lấy ngẫu nhiên $5$ thẻ từ $100$ thẻ nên số phần tử của không gian mẫu $  n\left(\Omega\right)=\mathrm{C}_{100}^5$.
			\itemch 
			Gọi $A$ là biến cố \lq\lq $5$ thẻ  lấy ra đều mang số chẵn\rq\rq.\\
			Vì từ $1$ đến $100$ có $50$ số chẵn, suy ra số  cách chọn $5$ thẻ đều mang số chẵn là $n(A)=\mathrm{C}_{50}^5.$\\
			Vậy xác suất để $5$ thẻ lấy ra đều mang số chẵn  là
			 $$ \mathrm{P}(A)=\dfrac{n(A)}{n\left(\Omega\right)}
			=\dfrac{\mathrm{C}_{50}^5}{\mathrm{C}_{100}^5}\approx 0{,}028 \ne  0{,}5.$$
			\itemch 
			 Gọi $B$ là biến cố \lq\lq $5$ thẻ lấy ra có $2$ thẻ mang số chẵn và $3$ thẻ mang số lẻ\rq\rq.\\
			Suy ra  $n(B)=\mathrm{C}_{50}^2\cdot\mathrm{C}_{50}^3$.\\
			Vậy $\mathrm{P}(B)=\dfrac{n(B)}{n\left(\Omega\right)} 
			 =\dfrac{\mathrm{C}_{50}^2\cdot\mathrm{C}_{50}^3}{\mathrm{C}_{100}^5}\approx 0{,}32$.
			\itemch 
			Gọi $C$ là biến cố \lq\lq Ít nhất một số ghi trên $5$ thẻ  được chọn chia hết cho $3$\rq\rq.\\
			Suy ra $\overline{C}\colon$  \lq\lq Cả $5$ số trên $5$ thẻ được chọn đều không chia hết cho $3$\rq\rq.\\
			Vì từ $1$ đến $100$ có $67$ số không chia hết cho $3$, suy ra số cách chọn $5$  thẻ đều đều không chia hết cho $3$ là $ n\left(\overline{C}\right)=\mathrm{C}_{67}^5$.\\
			Vậy  $ \mathrm{P}(C)=\dfrac{n(C)}{n\left(\Omega\right)}=\dfrac{\mathrm{C}_{100}^5-\mathrm{C}_{67}^5}{\mathrm{C}_{100}^5}\approx 0{,}87\ne 0{,}78$.
		\end{itemchoice}
	} 
\end{ex}

% nguồn: MaHoaNH_DungSai_Lop10
\begin{ex}%[0D0H2-5]%[ma:10360287]
	Một hộp có $5$ viên bi xanh, $6$ viên bi đỏ và $7$ viên bi vàng. Xét phép thử chọn ngẫu nhiên $3$ viên  bi. Khi đó
	 \choiceTF
		{\True Số phần tử  của gian mẫu của phép thử là $816$}
		{Xác xuất để chọn  được $3$ viên bi đỏ là $\dfrac{1}{272}$}
		{\True Xác xuất để chọn được  $3$ viên bi có đủ $3$ màu là $\dfrac{35}{136}$}
		{\True Xác xuất chọn được  nhiều nhất $2$ viên bi xanh là $\dfrac{403}{408}$}
	\loigiai{ 
		\begin{itemchoice}
			\itemch 
			Vì chọn $2$ viên bi trong $5+6+7=18$ viên bi nên số phần tử gian mẫu của phép thử là  $n\left(\Omega\right)=\mathrm{C}_{18}^3= 816$.
			 \itemch
			Gọi $A$ là biến  cố \lq\lq chọn được $3$ viên bi đỏ\rq\rq.\\
			Suy ra $n(A)=\mathrm{C}_{6}^3= 20$.\\ 
			Vậy xác suất để chọn  được $3$ viên bi đỏ là
			$\mathrm{P}(A)=\dfrac{n(A)}{n\left(\Omega\right)} 
			 =\dfrac{20}{816}=\dfrac{5}{204}\ne \dfrac{1}{272}$.
			 \itemch
			Gọi $B$ là biến cố  \lq\lq chọn được $3$ viên bi có đủ $3$ màu\rq\rq.\\
			Suy  ra $n(B)=\mathrm{C}_{5}^1\cdot\mathrm{C}_{6}^1\cdot\mathrm{C}_{7}^1= 210$.\\
			Vậy xác suất để chọn được $3$ viên có đủ $3$ màu là 
			$\mathrm{P}(B)=\dfrac{n(B)}{n\left(\Omega\right)} 
			=\dfrac{210}{816}=\dfrac{35}{216}$. 
			 \itemch
			Gọi  $C$ là biến cố \lq\lq chọn được nhiều nhất $2$ viên bi xanh\rq\rq.\\
			Suy ra $\overline{C}\colon$ \lq\lq chọn được  $3$ viên bi xanh\rq\rq.\\
			 Suy ra $n\left(\overline{C}\right)=\mathrm{C}_{5}^3=10$.
			Vậy $  \mathrm{P}(C)=1-\mathrm{P}\left(\overline{C}\right)=1-\dfrac{10}{816}=\dfrac{403}{408}$.
		\end{itemchoice}
	} 
\end{ex}

% nguồn: MaHoaNH_DungSai_Lop10
\begin{ex}%[0D0H2-5]%[ma:5f17c006]
	Một bình đựng $16$ viên bi, trong đó có  $7$ viên bi trắng, $6$ viên bi đen và $3$ viên bi đỏ. Lấy ngẫu nhiên $4$ viên bi. Khi đó
	 \choiceTF
		{Số phần tử của không gian  mẫu là $\mathrm{A}_{16}^4$}
		{\True Xác suất lấy được đúng $4$ bi trắng là  $\dfrac{1}{52}$}
		{\True Xác suất lấy được đủ $3$ màu là $\dfrac{9}{20}$} 
		{Xác  suất lấy được đúng $2$ màu $\dfrac{11}{20}$}
	\loigiai{ 
		\begin{itemchoice}
			 \itemch
			Chọn ra ngẫu nhiên $4$ viên bi trong $16$.  Số phần tử không gian mẫu là $n\left(\Omega\right)=\mathrm{C}_{16}^4$.
			 \itemch
			Gọi $A$ là biến  cố \lq\lq Lấy được đúng $4$ viên bi trắng\rq\rq.\\
			Số kết quả thuận lợi của biến cố  $A$ là $\mathrm{C}_7^4$.\\
			Vậy $\mathrm{P}(A)=\dfrac{\mathrm{C}_7^4}{\mathrm{C}_{16}^4}=\dfrac{1}{52}$. 
			 \itemch
			Gọi $B$ là biến cố \lq\lq Lấy được đủ $3$  màu\rq\rq.\\
			Lấy  được $1$ bi màu trắng, $1$ bi màu đen và 2 bi màu đỏ có $\mathrm{C}_7^1\cdot\mathrm{C}_6^1\cdot \mathrm{C}_3^2$ cách.\\
			Lấy được $1$ bi màu trắng, $2$ bi màu đen  và $1$ bi màu đỏ có $\mathrm{C}_7^1\cdot \mathrm{C}\mathrm{C}_6^2\cdot\mathrm{C}_3^1$ cách.\\
			Lấy được  $2$ bi màu trắng, $1$ bi màu đen và $1$ bi màu đỏ có $\mathrm{C}_7^2\cdot\mathrm{C}_6^1\cdot\mathrm{C}_3^1$ cách.\\
			Do  đó $n(B)=\mathrm{C}_7^1\cdot\mathrm{C}_6^1\cdot\mathrm{C}_3^2+\mathrm{C}_7^1\cdot\mathrm{C}_6^2\cdot\mathrm{C}_3^1+\mathrm{C}_7^2\cdot\mathrm{C}_6^1\cdot\mathrm{C}_3^1=819$.\\
			Suy  ra $\mathrm{P}(B)=\dfrac{819}{\mathrm{C}_{16}^4}=\dfrac{9}{20}$.
			\itemch 
			Gọi $C$ là biến  cố \lq\lq Lấy được đúng $2$ màu\rq\rq.\\
			Lấy được  đúng $1$ màu có $\mathrm{C}_7^4+\mathrm{C}_6^4$ cách.\\
			Lấy được đủ $3$ màu có  $\mathrm{C}_7^1\cdot\mathrm{C}_6^1\cdot\mathrm{C}_3^2+\mathrm{C}_7^1\cdot\mathrm{C}_6^2\cdot\mathrm{C}_3^1+\mathrm{C}_7^2\cdot\mathrm{C}_6^1\cdot\mathrm{C}_3^1$ cách.\\
			Do  đó $n(C)=\mathrm{C}_{16}^4-\left(\mathrm{C}_7^4+\mathrm{C}_6^4+
			\mathrm{C}_7^1\cdot\mathrm{C}_6^1\cdot\mathrm{C}_3^2+\mathrm{C}_7^1\cdot\mathrm{C}_6^2\cdot\mathrm{C}_3^1+\mathrm{C}_7^2\cdot\mathrm{C}_6^1\cdot\mathrm{C}_3^1\right)=951$ .\\ 
			 Vậy $P(B)=\dfrac{951}{\mathrm{C}_{16}^4}=\dfrac{951}{1820}$.
		\end{itemchoice}
	 }
\end{ex}

% nguồn: MaHoaNH_DungSai_Lop10
\begin{ex}%[0D0V2-5]%[ma:c1a3d326]
	Một  hộp chứa ba quả cầu trắng và bốn quả cầu đen. Lấy ngẫu nhiên đồng thời ba quả cầu. Các mệnh đề sau đây là đúng hay sai?
	 \choiceTF
		{\True Xác suất để  lấy được ba quả cầu đen là $\dfrac{4}{35}$}
		{\True Xác suất để lấy được một quả  cầu trắng và hai quả cầu đen là $\dfrac{18}{35}$}
		 {Xác suất để lấy được ba quả cầu cùng màu là $\dfrac{2}{7}$}
		{\True Xác suất để lấy được ít nhất một quả cầu màu trắng  là $\dfrac{31}{35}$}
	\loigiai{ 
		Số phần tử của không gian mẫu  $n(\Omega)=C^3_7=35$.
		\begin{itemchoice}
			\itemch  Gọi $A$ là biến cố: ``Lấy được ba quả cầu đen''. Khi đó $n(A)=C^3_4=4 \Rightarrow P(A)=\dfrac{n(A)}{n(\Omega)}=\dfrac{4}{35}$.
			\itemch Gọi $B$ là biến cố: ``Lấy được  một quả cầu trắng và hai quả cầu đen''. \\Khi đó $n(B)=C^1_3\cdot C^2_4=18 \Rightarrow P(B)=\dfrac{n(B)}{n(\Omega)}=\dfrac{18}{35}$.
			\itemch Gọi $C$ là biến cố: ``Lấy được ba  quả cầu cùng màu''.\\ Khi đó $n(C)=C^3_3+C^3_4=5 \Rightarrow P(C)=\dfrac{n(C)}{n(\Omega)}=\dfrac{5}{35}=\dfrac{1}{7}$.
			\itemch Gọi $D$ là biến cố: ``Lấy được ít nhất một quả cầu trắng''.\\ Suy ra $\overline{D}$  là biến cố: ``Lấy được ba quả cầu đen''. Suy ra $\overline{D}=A$.\\ Khi đó $P(D)=1-P(A)= \dfrac{31}{35}$.
		\end{itemchoice}
	} 
\end{ex}

% nguồn: MaHoaNH_DungSai_Lop10
\begin{ex}%[0D0V2-5]%[ma:829a39da]
	Có $100$ tấm thẻ được đánh  số từ $1$ đến $100$. Lấy ngẫu nhiên $5$ thẻ từ $100$ tấm thẻ ban đầu. Khi đó
	 \choiceTF
		{\True Số phần tử của không gian  mẫu là $\mathrm{C}_{100}^5$}
		{\True Xác suất của biến cố $A$: \lq\lq $5$ thẻ được chọn đều mang số chẵn\rq\rq\ là $\dfrac{1\,081}{38\,412}$} 
		{\True Xác suất của biến cố $B$:  \lq\lq $5$ thẻ được chọn có đúng $2$ thẻ mang số chẵn và $3$ thẻ mang số lẻ\rq\rq\ là $\dfrac{6\,125}{19\,206}$}
		{Xác suất của biến cố $C$: \lq\lq $5$ thẻ được chọn có cả thẻ mang số chẵn và thẻ mang  số lẻ\rq\rq\ là $0{,}9$}
	\loigiai{ 
		\begin{itemchoice}
			\itemch Chọn ngẫu nhiên $5$ thẻ  từ $100$ tấm thẻ ban đầu nên số phần tử không gian mẫu là $\mathrm{C}_{100}^5$.
			\itemch Chọn ngẫu nhiên $5$ thẻ từ $100$ tấm thẻ ban đầu trong đó $5$ thẻ được chọn đều mang số chẵn nên có số cách chọn  là $\mathrm{C}_{50}^5$ (cách chọn), khi đó $n(A)=\mathrm{C}_{50}^5$.\\
			Vậy  xác suất của biến cố $A$ là $\mathrm{P}(A)=\dfrac{n(A)}{n(\Omega)}=\dfrac{\mathrm{C}_{50}^5}{\mathrm{C}_{100}^5}=\dfrac{1\,081}{38\,412}$.
			\itemch Chọn ngẫu nhiên $5$ thẻ từ $100$ tấm thẻ ban đầu trong đó $5$ thẻ được chọn có $2$ thẻ mang số chẵn và $3$ thẻ mang số lẻ nên có số cách chọn là $\mathrm{C}_{50}^2\cdot \mathrm{C}_{50}^3$ (cách chọn).\\ 
			Do đó $  n(B)=\mathrm{C}_{50}^2\cdot \mathrm{C}_{50}^3$.\\
			Vậy xác suất của biến  cố $B$ là $\mathrm{P}(B)=\dfrac{n(B)}{n(\Omega)}=\dfrac{6\,125}{19\,206}$.
			\itemch Chọn ngẫu nhiên $5$ thẻ từ $100$ tấm thẻ ban đầu trong đó $5$ thẻ được chọn có cả thẻ mang số  chẵn và thẻ mang số lẻ nên nên có số cách chọn là $2\left(\mathrm{C}_{50}^1\cdot \mathrm{C}_{50}^4+\mathrm{C}_{50}^2\cdot \mathrm{C}_{50}^3\right)$ (cách chọn).\\
			Do đó $n(C)=2\left(\mathrm{C}_{50}^1\cdot \mathrm{C}_{50}^4+\mathrm{C}_{50}^2\cdot  \mathrm{C}_{50}^3\right)$.\\
			Vậy $\mathrm{P}(C)=\dfrac{n(C)}{n(\Omega)}=\dfrac{2\left(\mathrm{C}_{50}^1\cdot \mathrm{C}_{50}^4+\mathrm{C}_{50}^2\cdot \mathrm{C}_{50}^3\right)}{\mathrm{C}_{100}^{50}} \approx 0{,}944$. 
		\end{itemchoice}
	} 
\end{ex}

% nguồn: MaHoaNH_DungSai_Lop10
\begin{ex}%[0D0V2-5]%[ma:66141beb]
	Xét phép thử: Gieo đồng thời hai con súc sắc cân đối đồng chất. Xác  định tính Đúng - Sai của các khẳng định sau.
	 \choiceTF
		{\True $n(\Omega)=36$} 
		{Xác suất của biến cố ``Tổng số chấm  thu được từ hai con súc sắc bằng 6'' là $\dfrac{5}{26}$}
		 {\True Xác suất của biến cố ``Hiệu số chấm thu được từ hai con súc sắc bằng 2'' là $\dfrac{2}{9}$}
		 {\True Xác suất của biến cố ``Tích số chấm trên hai con súc sắc là một số chính phương'' là $\dfrac{2}{9}$}
	\loigiai{ 
		\begin{itemchoice}
			\itemch Số phần tử không gian  mẫu là $n(\Omega)=6\times 6=36$.
			\itemch Gọi biến cố $A$: \lq\lq Tổng  số chấm thu được từ hai con súc sắc bằng 6\rq\rq.\\
			Ta có $A=\{(1;5),(2;4),(3;3),(5;1),(4;2)\}  \Rightarrow n(A)=5$.\\
			Do vậy $\mathrm{P}(A)=\dfrac{n(A)}{n(\Omega)}=\dfrac{5}{36}$. 
			\itemch Gọi biến cố $B$: \lq\lq Hiệu số chấm thu được từ hai con súc sắc bằng  2\rq\rq.\\
			Ta có $B=\{(1;3),(2;4),(3;5),(4;6),(3;1),(4;2),(5;3),(6;4)\}$. 
			Suy ra  $n(B)=8$.\\
			 Khi đó $\mathrm{P}(B)=\dfrac{n(B)}{n(\Omega)}=\dfrac{8}{36}=\dfrac{2}{9}$.
			\itemch Gọi biến cố $C$: \lq\lq Tích số chấm trên hai con súc  sắc là một số chính phương\rq\rq.\\
			Ta có $C=\{(1;1),(2;2),(3;3),(4;4),(5;5),(6;6),(1;4),(4;1)\} \Rightarrow n(C)=8$.\\ 
			Vậy $\mathrm{P}(C)=\dfrac{n(C)}{n(\Omega)}=\dfrac{8}{36}=\dfrac{2}{9}$. 
		\end{itemchoice}
	} 
\end{ex}

% nguồn: MaHoaNH_DungSai_Lop10
\begin{ex}%[0D0V2-5]%[ma:0e89b73f]
	Bộ bài tú lơ khơ có $52$ quân bài. Rút ngẫu nhiên ra $4$  quân bài. Xét tính đúng sai của các mệnh đề sau:
	 \choiceTF
		{ Xác suất của biến cố  $A$: \lq\lq Rút ra được tứ quý Át \rq\rq\ là $\dfrac{1}{52}$}
		{\True Xác suất của biến cố $B$: \lq\lq Rút ra được hai quân Át, hai quân K \rq\rq\ là  $\dfrac{36}{270\,725}$}
		{ Xác suất của biến cố $C$: \lq\lq Rút ra được ít nhất một quân Át  \rq\rq\ là $\dfrac{38\,916}{54\,145}$}
		{\True Xác suất của biến cố $D$: \lq\lq Rút ra được $4$ quân trong đó có đúng $2$ quân ở cùng một tứ quý và $2$ quân còn lại ở hai tứ quý  khác nhau\rq\rq\ là $\dfrac{82\,368}{270\,725}$}
	 \loigiai{
		\begin{itemchoice}
			 \itemch \textbf{\textbf{Sai.}}\\
			Ta có  số cách chọn ngẫu nhiên 4 quân bài là $\mathrm{C}_{52}^4=270\,725$.\\
			Suy ra số phần tử  của không gian mẫu là $n( \Omega)=270\,725$.\\
			Vì bộ bài chỉ có 1 tứ quý Át nên số  phần tử của biến cố $A$ là $n( A)=1$.\\
			Vậy xác suất của biến  cố $A$ là $\mathrm{P}( A)=\dfrac{n( A)}{n( \Omega)}=\dfrac{1}{270\,725}$.
			\itemch \textbf{Đúng.}\\ 
			Có $\mathrm{C}_{4}^2$ cách rút được hai quân Át và có $\mathrm{C}_{4}^2$ cách  rút được hai quân $K$ nên số phần tử của biến cố $B$ là $n( B)=\mathrm{C}_{4}^2\cdot \mathrm{C}_{4}^2=36$.\\
			Vậy xác suất của biến cố $B$  là $\mathrm{P}( B)=\dfrac{n( B)}{n( \Omega)}=\dfrac{36}{270\,725}$.
			\itemch \textbf{Sai.}\\ 
			Xét biến cố $\overline{C}$: \lq\lq Rút không được quân Át nào \rq\rq. Suy ra  $n( \overline{C})=\mathrm{C}_{48}^4=194\,580$.\\
			Vậy xác suất của biến cố $C$ là $\mathrm{P}( C)=1-\mathrm{P}( \overline{C})=1-\dfrac{n( \overline{C})}{n(  \Omega)}=1-\dfrac{194\,580}{270\,725}=\dfrac{15\,229}{54\,145}$.
			\itemch \textbf{Đúng.}\\ 
			Có $\mathrm{C}_{13}^{1}$ cách chọn ra  $1$ tứ quý. Ứng với tứ quý này có $\mathrm{C}_{4}^2$ cách chọn ra $2$ quân bài.\\
			Có $\mathrm{C}_{12}^2$ cách chọn ra $2$ tứ quý  từ $12$ tứ quý còn lại. Mỗi tứ quý này có $\mathrm{C}_{4}^{1}$ cách chọn ra $1$ quân bài.\\
			Suy ra  $n( D)=\mathrm{C}_{13}^{1} \cdot \mathrm{C}_{4}^2\cdot \mathrm{C}_{12}^2 \cdot (\mathrm{C}_{4}^{1})^2=82\,368$.\\
			Vậy xác suất  của biến cố $D$ là $\mathrm{P}( D)=\dfrac{n( D)}{n( \Omega)}=\dfrac{82\,368}{270\,725}$.
		\end{itemchoice}
	 }
\end{ex}

% nguồn: MaHoaNH_DungSai_Lop10
\begin{ex}%[0D0V2-5]%[ma:97d51f5b]
	Một hộp có $15$ quả cầu trắng, $5$ quả  cầu đen. Xét phép thử chọn ngẫu nhiên $3$ quả cầu.
	 \choiceTF
		{\True Không gian mẫu của phép thử là: $1\,140$} 
		{ Xác suất để chọn được $2$ quả cầu trắng là: $\dfrac{7}{76}$} 
		{\True Xác suất để chọn được ít nhất một quả  cầu đen là: $\dfrac{137}{228}$}
		{ Xác suất để chọn được $3$ quả cầu thuộc hai  loại khác nhau là: $\dfrac{35}{76}$}
	 \loigiai{
		\begin{itemchoice}
			\itemch \textbf{Đúng.}\\ 
			Không gian mẫu của phép thử $n(  \Omega)=\mathrm{C}_{20}^3=1\,140$.
			\itemch \textbf{Sai.}\\ 
			Gọi $A$ là biến cố \lq\lq Chọn được hai quả cầu trắng \rq\rq. Khi đó  ta cần chọn $2$ quả trắng, $1$ quả đen.\\
			 Suy ra $n( A)=\mathrm{C}_{15}^2\cdot \mathrm{C}_{5}^{1}=525$.\\
			Vậy xác suất của biến cố $A$ là $\mathrm{P}(  A)=\dfrac{n( A)}{n( \Omega)}=\dfrac{525}{1\,140}=\dfrac{35}{76}$.
			\itemch \textbf{Đúng.}\\ 
			Gọi $B$ là biến cố \lq\lq Chọn được ít nhất một  quả cầu đen \rq\rq\ và $\overline{B}$ là biến cố \lq\lq Không chọn được quả đen nào \rq\rq, tức là chọn được $3$ quả trắng.\\
			 Suy ra $n( \overline{B})=\mathrm{C}_{15}^3=455$.\\
			Vậy xác  suất của biến cố $B$ là $\mathrm{P}( B)=1-\mathrm{P}( \overline{B})=1-\dfrac{455}{1\,140}=\dfrac{137}{228}$.
			\itemch  \textbf{Sai.}\\
			Gọi $C$ là biến cố \lq\lq  Chọn được ba quả cầu thuộc hai loại khác nhau\rq\rq.\\
			\textbf{Trường hợp 1}: Chọn $1$ quả trắng, $2$ quả đen  có $\mathrm{C}_{15}^{1}\cdot \mathrm{C}_{5}^2=150$ cách.\\
			\textbf{Trường hợp 2}: Chọn $2$  quả trắng, $1$ quả đen có $\mathrm{C}_{15}^2\cdot \mathrm{C}_{5}^{1}=525$ cách.\\
			 Suy ra $n( C)=150+525=675$.\\
			Xác suất  của biến cố $C$ là $\mathrm{P}( C)=\dfrac{n( C)}{n( \Omega)}=\dfrac{675}{1\,140}=\dfrac{45}{76}$.
		\end{itemchoice}
	 }
\end{ex}

% nguồn: MaHoaNH_DungSai_Lop10
\begin{ex}%[0D0V2-5]%[ma:0f22c26f]
	Từ một hộp chứa $12$ quả bóng gồm  $5$ quả màu đỏ và $7$ quả màu xanh, lấy ngẫu nhiên đồng thời $3$ quả. Khi đó
	 \choiceTF
		{\True Số  phần tử của không gian mẫu là $n\left(\Omega\right)=220$}
		{Xác suất để chọn được $3$ quả màu đỏ là $\dfrac{7}{44}$} 
		{\True Xác suất để chọn được $3$ quả trong đó có $2$ quả màu  đỏ, một quả màu xanh là $\dfrac{7}{22}$}
		{\True Xác suất để chọn được $3$ quả trong đó có ít nhất $1$ quả màu đỏ  là $\dfrac{37}{44}$}
	\loigiai{ 
		\begin{itemchoice}
			\itemch  Số phần tử của không gian mẫu là $n(\Omega)=\mathrm{C}_{12}^3=220$.
			\itemch Gọi $A$ là biến cố: \lq\lq Chọn được $3$  quả màu đỏ\rq\rq.\\
			Chọn $3$ quả màu đỏ trong $5$ quả màu  đỏ có $\mathrm{C}_5^3$ cách.\\
			 Do đó $n(A)=\mathrm{C}_5^3=10$.\\
			Vậy $\mathrm{P}(A)=\dfrac{10}{220}=\dfrac{1}{22}$. 
			\itemch Gọi $B$ là biến cố: \lq\lq Chọn được $3$ quả trong  đó có $2$ quả màu đỏ, một quả màu xanh\rq\rq.\\
			Chọn được $3$ quả trong đó có  $2$ quả màu đỏ, một quả màu xanh có $\mathrm{C}_5^2\cdot \mathrm{C}_7^1$ cách.\\
			 Do đó $n(B)=\mathrm{C}_5^2\cdot \mathrm{C}_7^1=70$.\\
			Vậy $\mathrm{P}(B)=\dfrac{70}{220}=\dfrac{7}{22}$. 
			\itemch Gọi $C$ là biến cố:  \lq\lq Chọn được $3$ quả trong đó có ít nhất $1$ quả màu đỏ\rq\rq.
			\begin{itemize}
				\item \textbf{Trường hợp 1.} Chọn $3$ quả gồm $1$ quả màu đỏ; $2$ quả màu xanh thì có $\mathrm{C}_5^1\cdot  \mathrm{C}_7^2$ cách chọn.
				\item \textbf{Trường hợp 2.}  Chọn $3$ quả gồm $2$ quả màu đỏ; $1$ quả màu xanh thì có $\mathrm{C}_5^2\cdot \mathrm{C}_7^1$ cách chọn.
				\item \textbf{Trường hợp 3.} Chọn $3$ quả màu  đỏ thì có $\mathrm{C}_5^3$ cách chọn.
			\end{itemize}
			Vậy  theo quy tắc cộng, ta có $n(C)=\mathrm{C}_5^1\cdot \mathrm{C}_7^2+\mathrm{C}_5^2\cdot \mathrm{C}_7^1+\mathrm{C}_5^3=185$.\\
			Vậy  $\mathrm{P}(A)=\dfrac{185}{220}=\dfrac{37}{44}$.
		\end{itemchoice}
	} 
\end{ex}

% nguồn: MaHoaNH_DungSai_Lop10
\begin{ex}%[0D0V2-5]%[ma:6822aeeb]
	Một hộp có $12$ viên bi, trong đó có $7$ viên bi xanh và $5$ viên  bi đỏ. Chọn ngẫu nhiên $5$ viên bi trong hộp.
	 \choiceTF
		{\True Số phần tử của  không gian mẫu là $792$}
		{\True Xác suất của biến  cố $A$: \lq\lq Năm viên bi đều là màu xanh\rq\rq\ là $\dfrac{7}{264}$}
		{ Xác suất của biến cố $B$: \lq\lq Trong $5$ viên bi lấy được có $3$ bi xanh  và $2$ bi đỏ\rq\rq\ là $\dfrac{125}{462}$}
		{ Xác suất của biến cố $C$: \lq\lq Trong 5 viên bi  lấy được có ít nhất 3 bi đỏ\rq\rq\ là $\dfrac{125}{396}$}
	\loigiai{ 
		\begin{itemchoice}
			 \itemch \textbf{Đúng.}\\
			Không gian mẫu là tập  tất cả các tập con gồm $5$ viên bi từ $12$ viên bi. Vậy số phần tử của không gian mẫu là $n( \Omega)=\mathrm{C}_{12}^5=792$.
			\itemch \textbf{Đúng.}\\ 
			Chọn  $5$ bi xanh từ $7$ bi xanh, có $\mathrm{C}_{7}^3=21$ (cách chọn). Suy ra $n(A)=21$.\\
			Vậy xác suất của biến cố  $A$ là $\mathrm{P}(A)=\dfrac{n(A)}{n( \Omega)} = \dfrac{21}{792} = \dfrac{7}{264}$.
			\itemch \textbf{Sai.}\\ 
			Mỗi phần tử của $B$  được hình thành từ 2 bước:\\
			\textbf{Bước 1}: Chọn $3$ viên bi xanh từ $7$ viên bi xanh, có $\mathrm{C}_{7}^3=35$  (cách chọn).\\
			\textbf{Bước 2}:  Chọn $2$ viên bi đỏ từ $5$ viên bi đỏ, có $\mathrm{C}_{5}^2=10$ (cách chọn).\\
			 Theo quy tắc nhân, ta có $n(B)=35\cdot 10=350$.\\
			Vậy $\mathrm{P}( B)=\dfrac{350}{792}=\dfrac{175}{396}$ 
			\itemch \textbf{Sai.}\\ 
			Trong $5$ viên bi lấy được có ít  nhất $3$ bi đỏ, có 3 trường hợp:\\
			\textbf{Trường hợp 1}: Trong $5$ viên bi  được chọn có $2$ bi xanh và $3$ bi đỏ có $\mathrm{C}_{7}^2\cdot \mathrm{C}_{5}^3=210$ (cách chọn)\\
			\textbf{Trường hợp 2}: Trong $5$ viên bi được chọn có $1$ bi xanh và $4$ bi  đỏ có $\mathrm{C}_{7}^{1}\cdot \mathrm{C}_{5}^4=35$ (cách chọn).\\
			\textbf{Trường hợp 3}: Trong $5$ viên bi  được chọn có 5 bi đỏ có $\mathrm{C}_{5}^5=1$ (cách chọn).\\
			Theo quy tắc cộng,  ta có $n(C)=210+35+1=246$.\\
			Vậy  $\mathrm{P}(C)=\dfrac{246}{792}=\dfrac{41}{132}$.
		\end{itemchoice}
	 }
\end{ex}

% nguồn: MaHoaNH_DungSai_Lop10
\begin{ex}%[0D0V2-5]%[ma:0423621c]
	Trong một chiếc  hộp có $20$ viên bi, trong đó có $8$ viên bi màu đỏ có bán kính khác nhau, $7$ viên bi màu xanh có bán kính khác nhau và $5$ viên bi màu vàng có bán kính khác nhau.
	Lấy ngẫu nhiên ra $3$  viên bi. Khi đó
	\choiceTF 
		{\True Số phần tử của không gian mẫu  là $n(\Omega)=1\,140$}
		{Xác suất để lấy được $3$ viên bi màu đỏ là  $\dfrac{14}{283}$}
		{Xác  suất lấy được $3$ viên bi không có viên bi nào màu đỏ là $\dfrac{11}{56}$}
		{\True Xác suất để lấy được $3$  viên bi có đúng hai màu $\dfrac{253}{380}$}
	 \loigiai{
		Gọi các biến cố 
		\begin{itemize}
			\item $A$: \lq\lq$3$ viên bi lấy ra đều  màu đỏ\rq\rq.
			\item $B$: \lq\lq$3$ viên bi không có viên bi  nào màu đỏ\rq\rq.
			\item  $C$: \lq\lq$3$ viên bi lấy ra có đúng hai màu\rq\rq.
		\end{itemize}
		\begin{itemchoice}
			\itemch Số cách lấy $3$ viên bi từ  $20$ viên bi là $\mathrm{C}_{20}^3$ nên ta có $n(\Omega)=\mathrm{C}_{20}^3=1\,140$.
			\itemch Số cách lấy $3$  viên bi màu đỏ là $n(A)=\mathrm{C}_8^3=56$.\\
			 Do đó $\mathrm{P}(A)=\dfrac{56}{1\,140}=\dfrac{14}{285}$.
			\itemch Số cách lấy $3$ viên bi không có viên bi  nào màu đỏ là $n(B)=\mathrm{C}_{12}^3=220$.\\
			Do đó xác suất lấy được $3$ bi không có bi nào màu đỏ  là $\mathrm{P}(B)=\dfrac{220}{1\,140}=\dfrac{11}{57}$.
			\itemch Số cách lấy $3$ viên bi  có đúng hai màu là
			\begin{itemize}
				\item Đỏ và  xanh: $\mathrm{C}_{15}^3-\left(\mathrm{C}_8^3+\mathrm{C}_7^3\right)$.
				\item  Đỏ và vàng: $\mathrm{C}_{13}^3-\left(\mathrm{C}_8^3+\mathrm{C}_5^3\right)$.
				 \item Vàng và xanh: $\mathrm{C}_{12}^3-\left(\mathrm{C}_5^3+\mathrm{C}_7^3\right)$.
			\end{itemize}
			Nên số cách lấy $3$ viên bi có đúng  hai màu là \[n(C)=\mathrm{C}_{15}^3+\mathrm{C}_{13}^3+\mathrm{C}_{12}^3-2\left(\mathrm{C}_8^3+\mathrm{C}_7^3+\mathrm{C}_5^3\right)=759.\]
			Vậy xác suất lấy được $3$ viên bi có đúng  $2$ màu là $\mathrm{P}(C)=\dfrac{759}{1\,140}=\dfrac{253}{380}$.
		\end{itemchoice}
	 }
\end{ex}

% nguồn: MaHoaNH_DungSai_Lop10
\begin{ex}%[0D0V2-5]%[ma:bd606a36]
	Trong một hộp có $50$ chiếc thẻ được đánh số từ $1$ đến  $50$. Rút ngẫu nhiên đồng thời $3$ chiếc thẻ từ hộp.
	\choiceTF 
		 {\True Số phần tử của không gian mẫu với phép thử trên là $n(\Omega)= 19600$}
		{\True Xác suất để rút được $3$ chiếc thẻ đều ghi số lẻ  bằng $\dfrac{23}{196}$}
		{Xác suất để rút được $3$ chiếc thẻ trong đó có  ít nhất một thẻ ghi số chẵn bằng $\dfrac{75}{196}$}
		{Xác suất để  tổng ba số trên ba thẻ rút được là số chia hết cho $3$ bằng $\dfrac{733}{2450}$}
	\loigiai{ 
	\begin{itemize}
		\item Số phần tử của không gian mẫu  với phép thử trên là $n(\Omega)= \mathrm{C}_{50}^3= 19600$ .
		\item Gọi $A$ là biến cố: \lq\lq Rút được $3$ chiếc thẻ đều  ghi số lẻ\rq\rq.\\
		Trong $50$ thẻ đã  cho có $25$ thẻ ghi số lẻ và $25$ thẻ ghi số chẵn.\\
		Do đó, $n(A)=\mathrm{C}_{25}^3= 2300$. Vậy $\mathrm{P}(A)=  \dfrac{n(A)}{n(\Omega)}= \dfrac{23}{196}$.
		\item Gọi $B$ là biến cố: \lq\lq Rút được $3$ chiếc thẻ trong đó có ít nhất một thẻ ghi  số chẵn\rq\rq.\\
		Ta thấy $B= \overline{A}$. Do đó $\mathrm{P}(B)=1-\mathrm{P(A)}= 1-\dfrac{23}{196}=  \dfrac{173}{196}$.\\
		\item Gọi $C$ là biến cố: \lq\lq Tổng ba số trên ba thẻ rút được là số chia  hết cho $3$\rq\rq.\\
		Trong $50$  thẻ đã cho ta chia thành 3 nhóm
		\begin{itemize}
			\item Nhóm 1: Các thẻ chia hết cho $3$, có $16$ thẻ. 
			\item Nhóm 2: Các thẻ chia cho $3$  dư $1$, có $17$ thẻ.
			\item Nhóm 2: Các thẻ  chia cho $3$ dư $2$, có $17$ thẻ.
		\end{itemize}
		Để rút được $3$ thẻ có tổng là số chia hết cho $3$,  có hai trường hợp xảy ra như sau
		\begin{itemize}
			 \item Trường hợp 1: Ba thẻ thuộc cùng một nhóm, có $\mathrm{C}_{16}^3+ \mathrm{C}_{17}^3+ \mathrm{C}_{17}^3= 1920$.
			\item Trường hợp 2: Trong ba thẻ, mỗi thẻ thuộc một nhóm  khác nhau, có $16\cdot 17 \cdot 17= 4624$ (cách).
		\end{itemize}
		Do đó $n(C)= 1920+4624=6544$. Vậy $\mathrm{P}(C)= \dfrac{n(C)}{n(\Omega)}= \dfrac{409}{1225}$. 
	\end{itemize}
	\begin{itemchoice}
		\itemch Đúng. 
		\itemch  Đúng.
		\itemch Sai. 
		 \itemch Sai.
	\end{itemchoice}
	 }
\end{ex}

% nguồn: MaHoaNH_DungSai_Lop10
\begin{ex}%[0D0V2-5]%[ma:2db37798]
	Trong một hộp có $40$ cái thẻ được đánh số từ $1$ đến $40$. Rút ngẫu nhiên đồng thời $3$ chiếc thẻ từ  hộp. Khi đó
	\choiceTF 
		{\True Số phần tử của không gian mẫu của phép thử trên là $  n\left(\Omega\right)=9880$}
		{\True Xác suất  để rút được $3$ chiếc thẻ đều ghi số lẻ bằng $\dfrac{3}{26}$}
		{Xác suất để rút được $3$ chiếc thẻ trong đó có  ít nhất một thẻ ghi số chẵn bằng $\dfrac{5}{13}$}
		{\True Xác suất để tổng ba số trên ba thẻ rút  được là số chia hết cho $3$ bằng $\dfrac{127}{380}$}
	 \loigiai{
		\begin{itemchoice}
			\itemch 
			 Rút ngẫu nhiên $3$ thẻ trong $40$ thẻ.\\
			Số phần  tử không gian mẫu là $n\left(\Omega\right)=\mathrm{C}_{40}^3=9\,880$.
			 \itemch
			Gọi $A$ là biến cố \lq\lq rút được $3$ chiếc thẻ đều ghi số lẻ\rq\rq.\\ 
			Trong $40$ thẻ được đánh số từ $1$ đến $40$ có $20$ thẻ đánh số lẻ và  $20$ thẻ đánh số chẵn.\\
			 Do đó số kết quả thuận lợi của biến cố $A$ là $n(A)=\mathrm{C}_{20}^3$.\\
			Vậy $\mathrm{P}(A)=\dfrac{\mathrm{C}_{20}^3}{\mathrm{C}_{40}^3}=\dfrac{3}{26}$. 
			\itemch 
			 Gọi $B$ là biến cố \lq\lq rút được $3$ chiếc thẻ trong đó có ít nhất một thẻ ghi số chẵn\rq\rq.\\
			Suy ra  $\overline{B}\colon$ \lq\lq rút được $3$ chiếc thẻ đều ghi số lẻ\rq\rq.\\
			Ta có $n\left(\overline{B}\right)=\mathrm{C}_{20}^3$.  \\
			Suy ra $\mathrm{P}(B)=1-\mathrm{P}\left(\overline{B}\right) 
			=1-\dfrac{\mathrm{C}_{20}^3}{\mathrm{C}_{40}^3}=\dfrac{23}{26}\ne \dfrac{5}{13}$. 
			\itemch 
			Gọi $C$ là biến cố \lq\lq tổng ba số trên ba thẻ rút được là số chia  hết cho $3$\rq\rq.\\
			 Trong $40$ thẻ được đánh số từ $1$ đến $40$ ta chia thành ba nhóm như sau:\\
			Nhóm 1: Các thẻ có số chia hết cho $3$ gồm các thẻ có số thuộc tập hợp $\{3, 6, 9,  \ldots, 39\}$, có $13$ thẻ.\\
			Nhóm 2: Các  thẻ có số chia cho $3$ dư $1$ gồm các thẻ có số thuộc tập hợp $\{1, 4, 7, \ldots, 40\}$, có $14$ thẻ.\\
			Nhóm 3: Các thẻ có số chia cho $3$ dư $2$ gồm các thẻ có số thuộc tập hợp $\{2,  5, 8, \ldots, 38\}$, có $13$ thẻ.\\
			Các trường  hợp chọn biến cố $D$ như sau
			\begin{itemize}
				\item {\bf Trường hợp 1:} Chọn $3$  thẻ cùng một nhóm.\\
				Số  cách chọn là $\mathrm{C}_{13}^3+\mathrm{C}_{14}^3+\mathrm{C}_{13}^3=936$.
				\item {\bf Trường hợp  2:} Chọn mỗi nhóm $1$ thẻ.\\
				Số cách chọn là  $\mathrm{C}_{13}^1\cdot\mathrm{C}_{14}^1\cdot\mathrm{C}_{13}^1=2\,366$.
			\end{itemize}
			 Do đó $n(C)=936+2\,366=3\,302$.\\
			Suy ra xác suất của biến cố  $C$ là
			$\mathrm{P}(C)=\dfrac{3\,302}{9\,880}=\dfrac{127}{380}$. 
		\end{itemchoice}
	} 
\end{ex}


%%%%%%%%%%%%%%%% TRẢ LỜI NGẮN %%%%%%%%%%%%%%%%

% nguồn: MaHoaNH_TraLoiNgan_Lop10
\begin{bt}%[0D0H2-5]%[ma:cf95da06]
	Trên giá sách có $4$ quyển  sách toán, $5$ quyển sách lý, $6$ quyển sách hóa. Lấy ngẫu nhiên $3$ quyển sách. Tính xác suất để $3$ quyển sách đươc lấy ra có ít nhất một quyển sách toán, kết quả làm tròn kết quả đến hàng phần trăm.
	 \shortans{0,64}
	\loigiai{Số phần tử của không  gian mẫu $n(\Omega)=\mathrm{C}_{15}^3$.\\
		Gọi $A$ là biến cố \lq\lq quyển sách được lấy ra có ít nhất một quyển sách  toán\rq\rq.\\
		 Ta có $n(A)=\mathrm{C}_{15}^3-\mathrm{C}_{11}^3$.\\
	Vậy xác suất cần  tìm là $$P(A)=\dfrac{n(A)}{n(\Omega)}=\dfrac{\mathrm{C}_{15}^3-\mathrm{C}_{11}^3}{\mathrm{C}_{15}^3}=\dfrac{58}{91}\approx0{,}64.$$}
\end{bt}

% nguồn: MaHoaNH_TraLoiNgan_Lop10
\begin{bt}%[0D0H2-5]%[ma:2ca6f646]
	Gieo  một con xúc xắc hai lần liên tiếp. Số khả năng thuận lợi để \lq\lq \,Tích số chấm ở hai lần gieo là một số nhỏ hơn $6$\,\rq\rq \ là bao nhiêu?\shortans{10}
	 \loigiai{
		 Số phần tử của không gian mẫu là $n\left(\Omega\right)=6\cdot 6=36$.\\
		Gọi $A$ là biến cố:  \lq\lq \,Tích số chấm ở hai lần gieo là một số nhỏ hơn $6$\,\rq\rq.\\
		Ta có $A=\left\{(1;1); (1;2); (2;1); (1;3); (3;1); (1;4); (4;1); (1;5);(5;1); (2;2)\right\}$.\\ 
		Số kết  quả thuận lợi cho biến cố $A$ là $n(A)=10$.\\
	 }
\end{bt}

% nguồn: MaHoaNH_TraLoiNgan_Lop10
\begin{bt}%[0D0H2-5]%[ma:80183d8d]
	Một lô hàng có $20$ sản phẩm trong đó có $4$ phế phẩm. Lấy  tùy ý $6$ sản phẩm từ lô hàng đó. Tính xác suất để $6$ sản phẩm được lấy có không quá $1$ phế phẩm? Làm tròn kết quả đến hàng phần trăm.
	\shortans{0,66} 
	\loigiai{ 
		 Số cách chọn có $1$ phế phẩm là $4\cdot \mathrm{C}^5_{16}$.\\
		Số  cách chọn không có phế phẩm là $\mathrm{C}^6_{16}$.\\
		Vậy  xác suất cần tìm là $\dfrac{4\cdot \mathrm{C}^5_{16}+\mathrm{C}^6_{16}}{\mathrm{C}^6_{20}}=\dfrac{637}{969}\approx 0{,}66$.
	} 
\end{bt}

% nguồn: MaHoaNH_TraLoiNgan_Lop10
\begin{bt}%[0D0H2-5]%[ma:9f506f9d]
	Trong một hộp chứa $4$ viên bi màu đỏ và $3$ viên bi màu xanh. Lấy ngẫu nhiên ra  $3$ viên bi. Hỏi có bao nhiêu cách lấy được $3$ viên cùng màu?
	\shortans{5} 
	 \loigiai{
		Để lấy được ba viên bi cùng màu, ta có hai  trường hợp
		\begin{itemize}
			\item Ba viên bi cùng màu xanh, có  $\mathrm{C}_3^3$ cách.
			\item Ba viên bi cùng  màu đỏ, có $\mathrm{C}_4^3$ cách.
		\end{itemize}
		Vậy số cách lấy ba viên bi cùng màu là $S=\mathrm{C}_3^3+\mathrm{C}_4^3=5$. 
	 }
\end{bt}

% nguồn: MaHoaNH_TraLoiNgan_Lop10
\begin{bt}%[0D0H2-5]%[ma:a990796d]
	Đề cương ôn tập môn Lịch sử có $30$ câu.  Đề thi được lập từ cách chọn ngẫu nhiên $10$ câu trong $30$ câu trong đề cương. Một học sinh chỉ học thuộc $25$ câu trong đề cương. Xác suất để trong đề thi có ít nhất $9$ câu hởi nằm trong $25$ câu mà học sinh đã học thuộc bằng bao nhiêu? (\textit{Kết quả làm tròn đến hàng phần trăm}).
	\shortans{0,45} 
	\loigiai{Không gian mẫu là số cách chọn $10$ câu trong $30$ câu.\\Suy ra số phần tử của không mẫu là $n(\Omega)=\mathrm{C}_{30}^{10}$. \\Gọi $A$ là biến cố \lq\lq Trong đề thi có ít nhất $9$ câu hỏi nằm  trong $25$ câu mà học sinh đã học thuộc \rq\rq.\\
		Để tìm số phần tử của $A$ ta xét các  trường hợp như sau:
		\begin{itemize}
			\item \textbf{Trường hợp 1.} $9$ câu thuộc và $1$ câu không thuộc:  có $\mathrm{C}_{25}^{9} \cdot \mathrm{C}_5^1$ khả năng.
			\item \textbf{Trường hợp 2.} $10$ câu đã  học thuộc hết: có $\mathrm{C}_{25}^{10}$ khả năng.
		\end{itemize}
		Suy ra số phần tử của  biến cố $A$ là $n(A)=\mathrm{C}_{25}^{9} \cdot \mathrm{C}_5^1+\mathrm{C}_{25}^{10}$.\\
		Vậy xác suất cần tính $\mathrm{P}(A)=\dfrac{\mathrm{C}_{25}^{9} \cdot  \mathrm{C}_5^1+\mathrm{C}_{25}^{10}}{\mathrm{C}_{30}^{10}}=\dfrac{3\,553}{7\,917} \approx 0{,}45$.
	} 
\end{bt}

% nguồn: MaHoaNH_TraLoiNgan_Lop10
\begin{bt}%[0D0H2-5]%[ma:6c16b572]
	Một bó hoa có $7$ bông hồng nhung, $6$ bông hồng bạch và $5$ bông hồng vàng. Rút ngẫu nhiên $7$ bông từ bó hoa. Tính xác suất để  rút được $3$ bông hồng nhung, $2$ bông hồng bạch và $2$ bông hồng vàng (\textit{Kết quả làm tròn đến chữ số thập phân thứ hai sau dấu phẩy}).
	\shortans{0,16} 
	 \loigiai{Số cách rút $7$ bông hoa từ $18$ bông hoa là $n(\Omega)=\mathrm{C}_{18}^{7}$ (cách).\\
		Gọi $A$ là biến cố \lq\lq Rút được $3$ bông hồng nhung, $2$ bông hồng bạch và $2$ bông hồng vàng \rq\rq.\\ 
		Một kết quả thuận lợi cho biến cố $A$ được thực  hiện theo các bước sau
		\begin{itemize}
			\item \textbf{Bước 1.} Rút $3$ bông hồng nhung có $\mathrm{C}_7^3$  cách.
			\item \textbf{Bước 2.}  Rút $2$ bông hồng bạch có $\mathrm{C}_{6}^2$ cách.
			\item \textbf{Bước 3.} Rút $2$ bông hồng vàng có  $\mathrm{C}_5^2$ cách.
		\end{itemize}
		Số kết quả thuận lợi cho biến cố  $A$ là $n(A)=\mathrm{C}_7^3 \cdot \mathrm{C}_{6}^2 \cdot \mathrm{C}_5^2$.\\
		Xác suất cần tìm là $\mathrm{P}(A)=\dfrac{\mathrm{C}_7^3 \cdot \mathrm{C}_{6}^2 \cdot \mathrm{C}_5^2}{\mathrm{C}_{18}^{7}}=\dfrac{875}{5\,304} \approx 0{,}16$. 
	} 
\end{bt}

% nguồn: MaHoaNH_TraLoiNgan_Lop10
\begin{bt}%[0D0H2-5]%[ma:92301e15]
	Một hộp  chứa $18$ quả bóng trong đó có $9$ quả màu trắng, $6$ quả màu đen và $3$ quả màu vàng. Lấy ngẫu nhiên đồng thời $5$ quả. Biết xác suất để trong $5$ quả bóng lấy ra có ít nhất $1$ quả màu vàng là $\dfrac{a}{b}$ (với $\dfrac{a}{b}$ là phân số tối giản). Tính $a+b$.
	\shortans{673} 
	\loigiai{ 
		Lấy $5$ quả bóng từ $18$ quả thì số phần tử của không gian mẫu $n(\Omega)=\mathrm{C}_{18}^{5}$.\\ 
		Gọi $A$  là biến cố \lq\lq Trong $5$ quả bóng lấy ra có ít nhất $1$ quả màu vàng\rq\rq.\\
		Khi đó $\overline{A}$ là biến cố \lq\lq Trong $5$ quả bóng lấy ra không có quả  nào màu vàng\rq\rq.\\
		Chọn  $5$ quả bóng không phải màu vàng từ $18-3=15$ quả có $\mathrm{C}_{15}^{5}$ cách.\\
		Suy ra số kết quả thuận lợi cho biến cố $\overline{A}$ là $n\left(\overline{A}\right)=\mathrm{C}_{15}^{5}$.\\ 
		 Xác suất của biến cố $\overline{A}$ là $\mathrm{P}\left(\overline{A}\right)=\dfrac{\mathrm{C}_{15}^{5}}{\mathrm{C}_{18}^{5}}=\dfrac{143}{408}$.\\
		Xác suất của biến cố $A$ là  $\mathrm{P}(A)=1-\mathrm{P}\left(\overline{A}\right)=1-\dfrac{143}{408}=\dfrac{265}{408}$.\\
		Xác suất để  trong $5$ quả bóng lấy ra có ít nhất $1$ quả màu vàng là $\dfrac{265}{408}$.\\
		Suy  ra $\heva{& a=265\\&b=408}\Rightarrow a+b=673$.
	 }
\end{bt}

% nguồn: MaHoaNH_TraLoiNgan_Lop10
\begin{bt}%[0D0H2-5]%[ma:98ba16da]
	Một hộp có $5$ viên bi đỏ, $3$ viên bi vàng và $4$ viên bi xanh. Chọn ngẫu nhiên từ hộp $4$ viên bi. Xác suất để $4$ viên bi được chọn có số bi đỏ lớn hơn số bi vàng và nhất thiết phải có  mặt bi xanh bằng $\dfrac{a}{b}$ với $\dfrac{a}{b}$ là phân số tối giản và $a$, $b \in \mathbb{Z}$. Tính $T=a+2 b$.
	 \shortans{82}
	\loigiai{ 
		Không gian mẫu là số cách chọn ngẫu nhiên $4$ viên bi từ hộp chứa $12$ viên bi. Suy ra số phần tử của không gian mẫu là $n(\Omega)=\mathrm{C}_{12}^4=495$.\\ 
		Gọi $A$ là  biến cố \lq\lq $4$ viên bi được chọn có số bi đỏ lớn hơn số bi vàng và nhất thiết phải có mặt bi xanh\rq\rq. Ta xét các trường hợp sau:
		\begin{itemize}
			\item \textbf{Trường hợp 1.} Chọn $1$ bi đỏ và $3$  bi xanh nên có $\mathrm{C}_5^1 \cdot \mathrm{C}_{4}^3$ cách.
			\item \textbf{Trường hợp 2.} Chọn $2$ bi đỏ  và $2$ bi xanh nên có $\mathrm{C}_5^2 \cdot \mathrm{C}_{4}^2$ cách.
			\item \textbf{Trường hợp 3.} Chọn $3$ bi đỏ và $1$  bi xanh nên có $\mathrm{C}_5^3 \cdot \mathrm{C}_{4}^1$ cách.
			\item \textbf{Trường hợp 4.} Chọn $2$ bi đỏ, $1$ bi vàng và $1$ bi xanh nên có $\mathrm{C}_5^2 \cdot \mathrm{C}_{3}^1 \cdot \mathrm{C}_{4}^1$ cách. 
		\end{itemize}
		Suy ra số phần tử của biến cố $A$ là $$n(A)=\mathrm{C}_5^1 \cdot \mathrm{C}_{4}^3+\mathrm{C}_5^2 \cdot \mathrm{C}_{4}^2+\mathrm{C}_5^3 \cdot \mathrm{C}_{4}^1+\mathrm{C}_5^2  \cdot \mathrm{C}_{3}^1 \cdot \mathrm{C}_{4}^1=240.$$
		Vậy xác suất cần tính $$\mathrm{P}(A)=\dfrac{n(A)}{n(\Omega)}=\dfrac{240}{495}=\dfrac{16}{33}\Rightarrow\heva{&a=16  \\& b=33} \Rightarrow T=a+2 b=82.$$
	} 
\end{bt}

% nguồn: MaHoaNH_TraLoiNgan_Lop10
\begin{bt}%[0D0H2-5]%[ma:2fcc650f]
	Trong giỏ có $6$ bông hoa  màu tím, $8$ bông hoa màu trắng. Lấy ngẫu nhiên $3$ bông hoa. Tính số khả năng thuận lợi để $3$ bông hoa lấy ra có đủ hai màu.
	\shortans{288} 
	\loigiai{ 
		Số phần tử của không  gian mẫu là $n\left(\Omega\right)=\mathrm{C}_{14}^3=364$.\\
		Gọi $A$ là biến cố $\colon$\lq\lq $3$ bông hoa  lấy ra có đủ hai màu\,\rq\rq.\\
		 Các trường hợp thuận lợi cho biến cố $A$ là
		\begin{itemize}
			\item[•]  Trường hợp $1$: Lấy được $1$ bông hoa tím và $2$ bông hoa trắng.\\
			Số  cách là $\mathrm{C}_6^1\cdot \mathrm{C}_8^2=168$ cách.
			\item[•] Trường hợp $2$: Lấy  được $2$ bông hoa tím và $1$ bông hoa trắng.\\
			Số cách là $\mathrm{C}_6^2\cdot \mathrm{C}_8^1=120$  cách.
		\end{itemize}
		Số kết quả thuận lợi cho $A$ là  $n(A)=168+120=288$.
	} 
\end{bt}

% nguồn: MaHoaNH_TraLoiNgan_Lop10
\begin{bt}%[0D0V2-5]%[ma:8122a448]
	Có $3$ viên bi đỏ và $7$ viên bi xanh, lấy ngẫu nhiên $4$ viên bi. Tính xác suất để lấy  được $2$ viên bi đỏ và $2$ viên bi xanh. \shortans{0,3}
	 \loigiai{
		 Xét phép thử lấy $4$ viên bi trong $3$ viên bi đỏ và $7$ viên bi xanh.\\
		Số phần tử của không gian mẫu là $n(\Omega)  = \mathrm{C}_{10}^{4} = 210$.\\
		Gọi $A\colon$ \lq\lq Lấy được $2$ viên bi đỏ và $2$ viên bi xanh \rq\rq.\\ 
		Khi đó $n(A) = \mathrm{C}_{3}^{2}\cdot \mathrm{C}_{7}^{2} = 63$.\\ 
		Xác suất để lấy được $2$ viên bi đỏ  và $2$ viên bi xanh là
		$$ \mathrm{P}(A) = \dfrac{n(A)}{n(\Omega)} = \dfrac{63}{210} =  \dfrac{3}{10}=0{,}3. $$
	} 
\end{bt}

% nguồn: MaHoaNH_TraLoiNgan_Lop10
\begin{bt}%[0D0V2-5]%[ma:e91698d7]
	Một hộp đựng $15$ viên bi, trong đó có $7$ viên bi xanh và $8$ viên bi đỏ. Lấy ngẫu nhiên ba viên bi (không kể thứ tự ra khỏi hộp). Tính xác suất để trong ba viên bi lấy ra có ít nhất một viên bi đỏ (làm tròn kết quả đến chữ số  thập phân thứ hai).
	\shortans{0,92} 
	 \loigiai{
		Số phần tử của không gian mẫu  $n(\Omega)=\mathrm{C}_{15}^3=455$.\\
		Gọi $A$ là biến cố \lq\lq trong ba viên bi lấy  ra có ít nhất một viên bi đỏ \rq\rq.\\
		Suy ra $\overline{A}$ là biến cố \lq\lq trong ba viên bi lấy ra  không có bi đỏ nào\rq\rq.\\
		Chọn $3$ bi xanh trong $7$ bi xanh có $\mathrm{C}_7^3=35$  cách.\\
		$$\Rightarrow  n(\overline{A})=35\Rightarrow \mathrm{P}(\overline{A})=\dfrac{n(\overline{A})}{n(\Omega)}=\dfrac{1}{13}.$$
		Vậy $\mathrm{P}(A)=1-\mathrm{P}(\overline{A})=\dfrac{12}{13}  \approx 0{,}92$.
	 }
\end{bt}

% nguồn: MaHoaNH_TraLoiNgan_Lop10
\begin{bt}%[0D0V2-5]%[ma:1a4f46d7]
	Chọn ngẫu nhiên $ 4 $ viên bi từ một hộp chứa $ 5 $ viên bi xanh, $ 6 $ viên bi đỏ, $ 7 $  viên bi vàng. Tính xác suất để chọn được $ 4 $ viên bi có đủ cả $ 3 $ màu, kết quả làm tròn đến hàng phần trăm.
	\shortans{0,51} 
	 \loigiai{
		Gọi $ A $ là biến cố: \lq\lq Chọn được $ 4 $ viên bi có đủ cả $ 3  $ màu\rq\rq.\\
		Số  phần tử của không gian mẫu là $ n(\Omega)=\mathrm{C}_{18}^4= 3060$.\\
		{\bf Trường hợp 1.} Chọn $ 2 $ bi xanh, $ 1 $ bi  đỏ, $ 1 $ bi vàng có $ \mathrm{C}_{5}^2\cdot \mathrm{C}_{6}^1 \cdot \mathrm{C}_{7}^1 =420$ cách.\\
		{\bf  Trường hợp 2.} Chọn $ 1 $ bi xanh, $ 2 $ bi đỏ, $ 1 $ bi vàng có $ \mathrm{C}_{5}^1\cdot \mathrm{C}_{6}^2 \cdot \mathrm{C}_{7}^1 =525$ cách.\\
		{\bf Trường hợp 3.}  Chọn $ 1 $ bi xanh, $ 1 $ bi đỏ, $ 2 $ bi vàng có $ \mathrm{C}_{5}^1\cdot \mathrm{C}_{6}^1 \cdot \mathrm{C}_{7}^2 =630$ cách.\\
		 Số phần tử của biến cố $ A $ là $ n(\Omega)=420+525+630=1575 $.\\
		Vậy xác suất cần  tìm là $ \mathrm{P}(A)=\dfrac{1575}{3060}= \dfrac{35}{68}\approx 0{,}51$.
	} 
\end{bt}

% nguồn: MaHoaNH_TraLoiNgan_Lop10
\begin{bt}%[0D0V2-5]%[ma:b4bc624a]
	Trong trò chơi \lq\lq Vòng quay may mắn\rq\rq, người chơi sẽ quay hai bánh xe. Mũi tên ở bánh xe thứ nhất có thể dừng ở một trong hai vị trí: Loại xe $50$cc và loại xe $110$cc. Mũi tên ở bánh xe thứ hai có thể dừng ở một trong bốn vị trí: màu đen, màu trắng, màu đỏ và màu xanh. Vị trí của mũi tên trên hai bánh xe sẽ xác định người chơi nhận được loại xe  và màu xe. Tính xác suất để người chơi nhận được loại xe $110$cc có màu trắng hoặc màu xanh.
	\shortans{0,25} 
	 \loigiai{
		Số phần tử của không gian mẫu  là $n(\Omega)=2\cdot 4=8$.\\
		Gọi $A$ là biến cố \lq\lq người chơi nhận được loại xe $110$cc có màu  trắng hoặc màu xanh\rq\rq.
		\begin{itemize}
			\item Quay được loại xe  $110$cc, có $1$ cách.
			\item Quay được màu trắng  hoặc xanh, có $2$ cách.
		\end{itemize}
		Suy ra $n(A)=1\cdot 2=2$.  Xác suất $\mathrm{P}(A)=\dfrac{2}{8}=\dfrac{1}{4}=0{,}25$.
	 }
\end{bt}

% nguồn: MaHoaNH_TraLoiNgan_Lop10
\begin{bt}%[0D0V2-5]%[ma:059ea680]
	Đề kiểm tra $15$ phút có $10$ câu trắc nghiệm, mỗi câu có bốn phương án trả lời, trong đó có một phương án đúng, trả lời đúng được $1{,}0$ điểm. Một thí sinh làm cả  $10$ câu, mỗi câu chọn một phương án. Xác suất để thí sinh đó đạt từ $8{,}0$ trở lên là $\dfrac{a}{4^{10}}$. Giá trị của $a$ bằng bao nhiêu?
	\shortans{436} 
	\loigiai{Số phân tử không gian mẫu $n(\Omega)=4^{10}$.\\Gọi $A$ là  biến cố \lq\lq Thí sinh đạt từ $8{,}0$ trở lên\rq\rq.
		Ta có các trường hợp: 
		\begin{itemize}
			\item \textbf{Trường hợp 1.} Thí sinh đúng $8$ câu, sai $2$ câu  có $\mathrm{C}_{10}^8 \cdot 3^2=405$ (cách).
			\item \textbf{Trường hợp 2.} Thí sinh đúng $9$ câu, sai $1$ câu có $\mathrm{C}_{10}^9 \cdot 3^1=30$  (cách).
			 \item \textbf{Trường hợp 3.} Thí sinh đúng cả $10$ câu có $\mathrm{C}_{10}^{10}=1$ (cách).
		\end{itemize}
		 Do đó $n(A)=405+30+1=436$.\\
		Vậy xác suất của biến cố  $A$ là $\mathrm{P}=\dfrac{n(A)}{n(\Omega)}=\dfrac{436}{4^{10}}$ nên $a=436$.
	} 
\end{bt}

% nguồn: MaHoaNH_TraLoiNgan_Lop10
\begin{bt}%[0D0V2-5]%[ma:afbad689]
	Từ một bộ đề thi học sinh giỏi môn Toán lớp $12$ mà mỗi đề gồm $5$ câu được chọn từ $20$ câu hỏi mức độ nhận biết, $15$ câu hỏi mức độ thông hiểu, $10$ câu hỏi mức độ vận dụng thấp và $5$ câu hỏi mức độ vận dụng cao. Lấy ngẫu nhiên $1$ đề thi trong bộ đề trên. Tính xác suất để đề thi lấy ra có đủ các câu hỏi thuộc $4$ mức độ nhận biết, thông hiểu, vận dụng thấp, vận dụng cao, đồng thời số câu mức độ vận dụng cao không quá $1$ (\textit{Kết  quả làm tròn đến chữ số thập phân thứ hai sau dấu phẩy}).
	 \shortans{0,15}
	 \loigiai{Số cách lấy ra ngẫu nhiên $1$ đề thi trong bộ đề là $n(\Omega)=\mathrm{C}_{50}^5$.\\
		Gọi $A$ là biến cố \lq\lq Đề thi lấy ra có đủ các câu hỏi thuộc $4$ mức độ nhận biết,  thông hiểu, vận dụng thấp, vận dụng cao, đồng thời số câu mức độ vận dụng cao không quá $1$ \rq\rq.
		\begin{itemize}
			\item \textbf{Trường hợp 1.} Đề thi có $5$ câu gồm $1$ câu mức độ vận dụng cao, $2$ câu mức độ vận dụng  thấp, $1$ câu mức độ thông hiểu và $1$ câu mức độ nhận biết. \\Có $\mathrm{C}_5^1 \cdot \mathrm{C}_{10}^2 \cdot \mathrm{C}_{15}^1 \cdot \mathrm{C}_{20}^1$ cách.
			\item \textbf{Trường hợp 2.} Đề thi có $5$ câu gồm $1$ câu mức độ vận dụng cao, $1$ câu mức độ vận dụng thấp, $2$ câu mức độ thông hiểu và $1$ câu mức  độ nhận biết. \\Có $\mathrm{C}_5^1 \cdot \mathrm{C}_{10}^1 \cdot \mathrm{C}_{15}^2 \cdot \mathrm{C}_{20}^1$ cách.
			\item \textbf{Trường hợp 3.} Đề thi có $5$ câu gồm $1$ câu mức độ vận dụng cao, $1$ câu mức độ vận dụng thấp, $1$ câu mức độ thông hiểu và $2$ câu mức độ nhận biết. \\Có $\mathrm{C}_5^1 \cdot \mathrm{C}_{10}^1 \cdot \mathrm{C}_{15}^1 \cdot  \mathrm{C}_{20}^2$ cách.
		\end{itemize}
		Suy ra $n(A)=\mathrm{C}_5^1 \cdot \mathrm{C}_{10}^2 \cdot \mathrm{C}_{15}^1 \cdot \mathrm{C}_{20}^1+\mathrm{C}_5^1 \cdot \mathrm{C}_{10}^1 \cdot \mathrm{C}_{15}^2 \cdot  \mathrm{C}_{20}^1+\mathrm{C}_5^1 \cdot \mathrm{C}_{10}^1 \cdot \mathrm{C}_{15}^1 \cdot \mathrm{C}_{20}^2=315\,000$.\\
		Vậy xác suất  cần tính là $\mathrm{P}(A)=\dfrac{n(A)}{n(\Omega)}=\dfrac{1\,125}{7\,567} \approx 0{,}15$.
	} 
\end{bt}

% nguồn: MaHoaNH_TraLoiNgan_Lop10
\begin{bt}%[0D0V2-5]%[ma:1de83b74]
	Một bộ đề thi toán học sinh giỏi lớp $12$ mà mỗi đề gồm $5$ câu được chọn từ $15$ câu dễ, $10$ câu trung bình và $5$ câu khó. Một đề thi được gọi là \lq\lq Tốt\rq\rq\, nếu trong đề  thi có cả ba câu dễ, trung bình và khó, đồng thời số câu dễ không ít hơn $2$. Lấy ngẫu nhiên một đề thi trong bộ đề trên. Tìm xác suất để đề thi lấy ra là một đề thi \lq\lq Tốt\rq\rq (\textit{Kết quả làm tròn đến chữ số thập phân thứ nhất sau dấu phẩy}).
	 \shortans{0,4}
	\loigiai{Số phần tử của không gian mẫu  là $n(\Omega)=\mathrm{C}_{30}^5=142\,506$.\\
		Gọi $A$ là biến  cố \lq\lq Đề thi lấy ra là một đề thi \lq\lq Tốt\rq\rq \rq.\\
		Vì trong một đề thi \lq\lq Tốt\rq\rq\, có cả ba câu dễ, trung bình và  khó, đồng thời số câu dễ không ít hơn $2$ nên ta có các trường hợp sau đây thuận lợi cho biến cố $A$ là
		\begin{itemize}
			\item \textbf{Trường hợp 1.} Đề thi gồm $3$ câu dễ, $1$ câu trung bình và $1$ câu khó, có $\mathrm{C}_{15}^3 \cdot \mathrm{C}_{10}^1 \cdot \mathrm{C}_5^1$  đề.
			\item \textbf{Trường hợp  2.} Đề thi gồm $2$ câu dễ, $2$ câu trung bình và $1$ câu khó, có $\mathrm{C}_{15}^2 \cdot \mathrm{C}_{10}^2 \cdot \mathrm{C}_5^1$ đề.
			\item \textbf{Trường hợp 3.} Đề thi gồm $2$ câu  dễ, $1$ câu trung bình và $2$ câu khó, có $\mathrm{C}_{15}^2 \cdot \mathrm{C}_{10}^1 \cdot \mathrm{C}_5^2$ đề.
		\end{itemize}
		Suy ra số phần tử của biến cố $A$ là $n(\Omega)=\mathrm{C}_{15}^3 \cdot \mathrm{C}_{10}^1 \cdot \mathrm{C}_5^1+\mathrm{C}_{15}^2  \cdot \mathrm{C}_{10}^2 \cdot \mathrm{C}_5^1+\mathrm{C}_{15}^2 \cdot \mathrm{C}_{10}^1 \cdot \mathrm{C}_5^2=56\,875$.\\
		Vậy xác suất cần tính  $\mathrm{P}(A)=\dfrac{n(\Omega_A)}{n(\Omega)}=\dfrac{56\,875}{142\,506}=\dfrac{625}{1\,566} \approx 0{,}4$.
	} 
\end{bt}

% nguồn: MaHoaNH_TraLoiNgan_Lop10
\begin{bt}%[0D0V2-5]%[ma:f5052462]
	Một hộp chứa $6$ quả bóng màu đỏ được đánh số từ $1$ đến $6$; $5$ quả bóng màu vàng được đánh số từ $1$ đến $5$ và $4$ quả bóng màu xanh được đánh số từ $1$ đến $4$. Lấy ngẫu nhiên $4$ quả bóng trong hộp. Tính xác suất để $4$ quả bóng lấy ra có đủ ba  màu đồng thời không có hai quả bóng nào được đánh số trùng nhau (\textit{Kết quả làm tròn đến hàng phần trăm}).
	 \shortans{0,16}
	\loigiai{Số phần tử của không gian  mẫu là số cách lấy bất kì $4$ quả bóng từ $15$ quả bóng.\\
		Suy ra số phần tử của không gian  mẫu là $n(\Omega)=\mathrm{C}_{15}^4=1\,365$.\\
		Gọi $A$ là biến  cố \lq\lq$4$ quả bóng lấy ra có đủ ba màu đồng thời không có hai quả bóng nào được đánh số trùng nhau\rq\rq.\\
		 Các trường hợp xảy ra biến cố $A$:
		\begin{itemize}
			\item \textbf{Trường hợp 1.} $4$ quả cầu lấy ra có $2$ xanh, $1$  vàng, $1$ đỏ, có $\mathrm{C}_{4}^2 \cdot \mathrm{C}_{3}^1 \cdot \mathrm{C}_{3}^1$ cách.
			\item \textbf{Trường hợp 2.} $4$ quả cầu lấy ra có $1$ xanh, $2$ vàng, $1$ đỏ, có $\mathrm{C}_{4}^1 \cdot \mathrm{C}_{4}^2 \cdot \mathrm{C}_{3}^1$ cách. 
			\item \textbf{Trường hợp 3.} $4$ quả cầu lấy ra có $1$ xanh, $1$  vàng, $2$ đỏ, có $\mathrm{C}_{4}^1 \cdot \mathrm{C}_{4}^1 \cdot \mathrm{C}_{4}^2$ cách.
		\end{itemize}
		Suy ra số  phần tử của biến cố $A$ là $$n(A)=\mathrm{C}_{4}^2 \cdot \mathrm{C}_{3}^1 \cdot \mathrm{C}_{3}^1+\mathrm{C}_{4}^1 \cdot \mathrm{C}_{4}^2 \cdot \mathrm{C}_{3}^1+\mathrm{C}_{4}^1 \cdot \mathrm{C}_{4}^1 \cdot \mathrm{C}_{4}^2=222.$$
		Do đó xác suất  của biến cố $A$ là $\mathrm{P}(A)=\dfrac{n(A)}{n(\Omega)}=\dfrac{222}{1\,365}=\dfrac{74}{455} \approx 0{,}16$.
	} 
\end{bt}

% nguồn: MaHoaNH_TraLoiNgan_Lop10
\begin{bt}%[0D0V2-5]%[ma:b4e37c4d]
	Một bài kiểm tra trắc nghiệm với $10$ câu hỏi, mỗi câu hỏi có $4$ đáp án và trong đó chỉ có duy nhất một đáp án là đúng. Xét phép thử $T$: Bạn Minh làm bài kiểm tra bằng cách chọn ngẫu nhiên các đáp án cho các câu hỏi. Biết mỗi câu trả lời đúng bạn được $1$ điểm và sai thì không mất điểm. Gọi $A$ là biến cố: \lq\lq Bạn Minh được ít nhất $5$ điểm\rq\rq. Tính xác suất của biến cố $A$ (làm tròn kết quả  đến hàng phần trăm).
	\shortans{0,08} 
	\loigiai{ 
		Vì mỗi câu hỏi có $4$ sự lựa chọn nên có $4^{10}$ cách làm bài kiểm tra với $10$ câu hỏi hay $\mathrm{n}(\Omega)=4^{10}$.\\ 
		\begin{itemize}
			\item TH1:  Bạn Minh được $5$ điểm có $\mathrm{C}_{10}^5\cdot3^5$ cách.
			\item TH2:  Bạn Minh được $6$ điểm có $\mathrm{C}_{10}^6\cdot3^4$ cách.
			\item  TH3: Bạn Minh được $7$ điểm có $\mathrm{C}_{10}^7\cdot3^3$ cách.
			\item TH4: Bạn Minh được $8$ điểm có $\mathrm{C}_{10}^8\cdot3^2$ cách. 
			\item TH5: Bạn Minh được $9$ điểm có $\mathrm{C}_{10}^9\cdot3$ cách. 
			\item  TH6: Bạn Minh được $10$ điểm có $\mathrm{C}_{10}^{10}$ cách.
		\end{itemize}
		Số phần tử của  biến cố $A$ là
		$$\mathrm{n}(A)=\mathrm{C}_{10}^6\cdot3^4+\mathrm{C}_{10}^7\cdot3^3+\mathrm{C}_{10}^8\cdot3^2+\mathrm{C}_{10}^9\cdot3+\mathrm{C}_{10}^{10}=81922.$$ 
		Xác suất của biến cố  $A$ là
		$$\mathrm{P}(A)=\dfrac{81922}{4^{10}}\approx0{,}08.$$ 
	 }
\end{bt}


%%%%%%%%%%%%%%%% TỰ LUẬN %%%%%%%%%%%%%%%%

% nguồn: MaHoaNH_TuLuan_Lop10
\begin{bt}%[0D0H2-5]%[ma:3fd406a9]
	Một hộp gồm $30$ quả cầu được đánh số từ  $1$ đến $30$. Chọn ngẫu nhiên $3$ quả cầu từ hộp đó. Tính xác suất để lấy được $3$ quả cầu có đúng $1$ quả cầu ghi số lẻ và tích $3$ số ghi trên ba quả cầu là một số chia hết cho $8$?
	\loigiai{Ta có số phần tử không gian mẫu là số cách chọn 3  quả cầu từ hộp $30$ quả cầu nên $n(\Omega)=\mathrm{C}_{30}^{3}$.\\
		Gọi $A$ là  biến cố thỏa mãn yêu cầu bài toán.\\
		Chọn  một số lẻ từ $30$ số có $15$ cách.\\
		Đặt $A_1=\left\{4;8;12;16;20;24;28\right\}$, $A_2=\left\{2;6;10;14;18;22;16;30\right\}$ trong đó $A_1$ gồm những số chia hết cho $4$ và $A_2$ là những số chẵn không chia  hết cho $4$.\\
		Do quả cầu đầu tiên mang số lẻ nên để chọn $3$ quả cầu thỏa mãn yêu  cầu bài toán thì tích $2$ số trên $2$ quả cầu còn lại phải là số chia hết cho $8$.\\
		\textbf{Trường hợp 1:} Chọn $2$  số từ tập $A_1$ có $\mathrm{C}_{7}^{2}=21$ cách.\\
		\textbf{Trường hợp 2:} Chọn $1$ số từ tập $A_1$ và số còn lại là số tùy  ý từ $A_2$ có $7 \cdot 8 =56$ cách.\\
		 Vậy $n(A)=15 \cdot (21+66)=1155$ nên $\mathrm{P}(A)=\dfrac{n(A)}{n(\Omega)}=\dfrac{1155}{\mathrm{C}_{30}^3}=\dfrac{33}{116}$.
	 }
\end{bt}

% nguồn: Ok_TuanLuan_Toan11_Lan1
\begin{bt}%[0D0H2-5]%[ma:507e7930]
	Trong một chiếc hộp có $15$ viên bi có cùng kích thước và khối lượng, trong đó có $4$ viên bi màu đỏ, $5$ viên bi màu xanh và $6$ viên bi màu vàng.  Lấy ngẫu nhiên đồng thời $3$ viên bi. Gọi $P$ là xác suất lấy được $3$ viên bi có đúng hai màu. Giá trị của $65\cdot P$ bằng bao nhiêu?
	 \loigiai{
		Chọn ngẫu nhiên đồng thời $3$ bi từ $15$ nên số phần tử không gian mẫu là  $n(\Omega)=\mathrm{C}_{15}^{3}=455$ (cách).\\
		\textbf{Trường hợp 1.} $3$  viên lấy ra cùng màu đỏ: $\mathrm{C}_{4}^{3}$.\\
		\textbf{Trường hợp  2.} $3$ viên lấy ra cùng màu xanh: $\mathrm{C}_{5}^{3}$.\\
		\textbf{Trường hợp  3.} $3$ viên lấy ra cùng màu vàng: $\mathrm{C}_{6}^{3}$.\\
		\textbf{Trường hợp 4.} $3$ viên lấy ra có đủ $3$ màu:  $4\cdot 5\cdot 6=120$.\\
		Số cách  lấy ra $3$ viên bi có đúng hai viên bi có đúng hai màu là $\mathrm{C}_{15}^{3}-\left(\mathrm{C}_{4}^{3}+\mathrm{C}_{5}^{3}+\mathrm{C}_{6}^{3}+120\right)=301$.\\
		Xác  suất để lấy được $3$ viên có đúng hai màu là $\mathrm{P}=\dfrac{301}{\mathrm{C}_{15}^{3}}=\dfrac{43}{65}$.\\
		Vậy $65\cdot \mathrm{P}=65\cdot  \dfrac{43}{65}=43$.
	} 
\end{bt}

% nguồn: Ok_TuanLuan_Toan11_Lan2
\begin{ex}%[0D0H2-5]%[ma:f9d03008]
	Trong một chiếc hộp có $15$ viên bi có cùng kích thước và khối lượng, trong đó có $4$ viên bi màu đỏ, $5$ viên  bi màu xanh và $6$ viên bi màu vàng. Lấy ngẫu nhiên đồng thời $3$ viên bi. Gọi $P$ là xác suất lấy được $3$ viên bi có đúng hai màu. Giá trị của $65\cdot P$ bằng bao nhiêu?
	 \loigiai{
		Chọn ngẫu nhiên đồng thời $3$ bi từ $15$ nên số phần tử không gian mẫu là $n(\Omega)=\mathrm{C}_{15}^{3}=455$ (cách).\\ 
		\textbf{Trường hợp 1.} $3$ viên lấy ra cùng màu đỏ: $\mathrm{C}_{4}^{3}$.\\ 
		\textbf{Trường hợp 2.}  $3$ viên lấy ra cùng màu xanh: $\mathrm{C}_{5}^{3}$.\\
		\textbf{Trường hợp 3.} $3$ viên lấy ra cùng màu vàng: $\mathrm{C}_{6}^{3}$.\\ 
		\textbf{Trường hợp 4.} $3$ viên lấy ra có đủ $3$ màu:  $4\cdot 5\cdot 6=120$.\\
		Số cách lấy  ra $3$ viên bi có đúng hai viên bi có đúng hai màu là $\mathrm{C}_{15}^{3}-\left(\mathrm{C}_{4}^{3}+\mathrm{C}_{5}^{3}+\mathrm{C}_{6}^{3}+120\right)=301$.\\
		Xác suất để lấy được $3$ viên  có đúng hai màu là $\mathrm{P}=\dfrac{301}{\mathrm{C}_{15}^{3}}=\dfrac{43}{65}$.\\
		Vậy $65\cdot  \mathrm{P}=65\cdot \dfrac{43}{65}=43$.
	} 
\end{ex}

% nguồn: MaHoaNH_TuLuan_Lop10
\begin{bt}%[0D0V2-5]%[ma:e9fcbd87]
	Trong tủ có $4$  đôi giày khác loại. Bạn Lan lấy ra ngẫu nhiên $2$ chiếc giày. Tính xác suất để lấy ra được một đôi giày hoàn chỉnh.
	\loigiai{ 
		Gọi $A$ là  biến cố \lq\lq Lấy ra được một đôi giày hoàn chỉnh\rq\rq.\\
		Số phần tử của không gian mẫu $n(\Omega)=\mathrm{C}_8^2=28$. \\ Số phần  tử của biến cố $A$ là $n(A)=4$.\\
		Vậy xác  suất của biến cố $A$ là $\mathrm{P}(A)=\dfrac{n(A)}{n(\Omega)}=\dfrac{4}{28}=\dfrac{1}{7}$.
	 }
\end{bt}

% nguồn: MaHoaNH_TuLuan_Lop10
\begin{bt}%[0D0V2-5]%[ma:6510df70]
	Một lớp học có $26$ bạn nam và $20$ bạn nữ. Chọn ngẫu nhiên một bạn trong lớp. Tính xác suất để  bạn được chọn là nam.
	 \loigiai{
		Số phần tử của  không gian mẫu là $n(\Omega)=26+20=46$.\\
		Gọi $A$ là biến cố  bạn được chọn là nam. Vì lớp học có $26$ bạn nam nên có $26$ cách chọn một bạn nam. Do đó, ta có $n(A)=26$.\\
		Vậy xác  suất của biến cố $A$ là $\mathrm{P}(A)=\dfrac{n(A)}{n(\Omega)}=\dfrac{26}{46}=\dfrac{13}{23}$.
	 }
\end{bt}
